% PLAN1-2026-11-11 preserved verbatim excerpts from 64-page main.tex
% SOURCE: growupkuriyama-hub/Papers, 02-v2-lg-fanout-boundary/main.tex
% SOURCE BLOB: d7bb4f7d3c51d1a050d9f082e03726f342543621
% FRAGMENTS ONLY: not an independently compilable paper, do not \input into plan1 main.tex.
% The original source and its citations remain in git history.

% [PART 1] Results on individual U1 typings and isolated letters formerly in Section 7
\paragraph{Individual typings onto $U_1$.}
Theorem~\ref{thm:fanout-one-data} concerns pseudovarieties.  For an
individual typing at fan-out one, what matters is which words realize a copy
$\{e,z\}$ of $U_1$: letter-disjoint words of type $e$ force exponential data,
while an isolated letter (Definition~\ref{def:isolated-letter}), whose type is
realized only by its own powers, gives polynomial data
(Theorem~\ref{thm:two-letter-fanout-one} below).  For typings onto $U_1$
these two cases are exhaustive (Corollary~\ref{cor:u1-classification}).

\begin{corollary}[A typing-level obstruction at fan-out one]
\label{cor:typing-level-obstruction}
Let $r_{\max}\ge2$ and let $s:\Sigma^+\twoheadrightarrow\mathsf S$ be a typing.
Suppose that distinct idempotents $e,z\in\mathsf S$ satisfy $ez=ze=z$, and that
words $u,v,w\in\Sigma^+$ satisfy $s(u)=s(v)=e$ and $s(w)=z$ and are
\emph{letter-disjoint}: no letter occurs twice in $uvw$.  Let $H_n$ be the image
of the language $L_n$ of Proposition~\ref{prop:three-letter-family}, over three
fresh letters, under $\mathtt a\mapsto u$, $\mathtt b\mapsto v$, $\mathtt c\mapsto w$,
and let $G''_n$ be the image of $G_n$, with rules $S\to E_n$, $E_0\to u$ and
$E_i\to vE_{i-1}E_{i-1}v\mid wE_{i-1}w$.  Then $H_n\in\Cmcf{1}{s}$, $G''_n$ is a
reduced good rank-two presentation of $H_n$ with
\[
 \|G''_n\|=(2|v|+2|w|+9)n+|u|+6,\qquad
 \bar\theta_{G''_n}\le2|u|+2|v|+4n|w|,
\]
and every characteristic sample of $H_n$ for $\mathcal B^{\le r_{\max}}_{1,s}$,
for $\mathcal B_{1,s}$, or for any set-driven operator having characteristic
samples throughout $\Cmcf{1}{s}$ has total encoded length greater than
$(|u|+2|v|)2^n-2|v|$.  So for this single typing no polynomial bounds the
characteristic-data length at fan-out one and rank $r_{\max}$ as in
Definition~\ref{def:uniform-polydata}.
\end{corollary}

\begin{proof}
$L_n$ is finite and Clark--Eyraud substitutable
(Proposition~\ref{prop:three-letter-family}), so $H_n$ is Clark--Eyraud
substitutable by Lemma~\ref{lem:app-u1-code}; it is finite, hence
$H_n\in\Cmcf{1}{s}$.  The grammar is good, and reduced since
$E_i\Rightarrow^*w^iuw^i$.  The size follows from
Definition~\ref{def:standard-mcfg}: the start rule has size $3$, the base rule
$1+|u|$, and the two rules of each level $2|v|+5$ and $2|w|+3$.  Completing every
other choice by the rules $E_j\to wE_{j-1}w$ shows that each rule has a
successful use of length at most $|u|+2n|w|$ (start and unary rules) or
$2|u|+2|v|+(2n+2i-4)|w|$ (the binary rule of level $i$).  In $\{e,z\}$, $e$ is an
identity and $z$ is absorbing, so a word of $H_n$ has type $e$ iff it is the
image of a word of $L_n$ without $\mathtt c$, that is, of $W_n$.  The image $h(W_n)$ has length
$(|u|+2|v|)2^n-2|v|$ by $|h(W_i)|=2|h(W_{i-1})|+2|v|$, and
Proposition~\ref{prop:rare-type} forces it into every such sample.
\end{proof}

With $u,v,w$ letters this is the case of two letters of type $e$; in
particular it applies to every typing onto $U_1$ with at least two letters of
type $1$, since surjectivity gives a letter of type $\zeta$.  It also applies
when no letter has type $e$: for $\mathsf S=\{1,g,z\}$ with identity $1$,
$g^2=1$ and zero $z$, and the typing $\mathtt a,\mathtt b,\mathtt c,\mathtt d\mapsto g$,
$\mathtt t\mapsto z$, take $u=\mathtt{ab}$, $v=\mathtt{cd}$, $w=\mathtt t$ and
$e=1$.  Whether two non-commuting words of type $e$ always suffice, without
letter-disjointness, is open in general; since two words commute iff they are
powers of one word, their absence means that all words of type $e$ are powers of
one primitive word.  For typings onto $U_1$ the question does not arise: with
two letters of type $1$ the corollary applies, and with one, $\mathtt a$ say, the
words of type $1$ are the powers of $\mathtt a$ and the data are polynomial
(Corollary~\ref{cor:u1-classification}).

Let $\Sigma=\{\mathtt a,\mathtt b\}$ and let $s_{U,2}:\Sigma^+\twoheadrightarrow U_1$ send $\mathtt a$ to
$1$ and $\mathtt b$ to $\zeta$ (the binary typing).  For this typing, with a single letter of type $1$,
Corollary~\ref{cor:typing-level-obstruction} does not apply; still, its data can
be quadratic in $\bar\theta_G$ for presentations of logarithmic size.

\begin{proposition}[Quadratic data for the binary typing at fan-out one]
\label{prop:binary-u1-quadratic}
For $m,n\ge1$ let $L_{m,n}\subseteq\{\mathtt a,\mathtt b\}^+$ be the set of all
concatenations of $2^n$ words from $\{\mathtt a^{2^m},\mathtt b\}$.
\begin{enumerate}[label=\textup{(\roman*)},leftmargin=*]
\item $L_{m,n}\in\Cmcf{1}{s_{U,2}}$.
\item $L_{m,n}$ has a reduced good fan-out-one presentation $G_{m,n}$ of rule rank
two with $\|G_{m,n}\|=6m+6n+9$ and $\bar\theta_{G_{m,n}}=2^m+2^n-1$.
\item For $\mathcal B_{1,s_{U,2}}$, for $\mathcal B^{\le r}_{1,s_{U,2}}$ \textup($r\ge1$\textup),
and for every set-driven operator with characteristic samples throughout
$\Cmcf{1}{s_{U,2}}$, every characteristic sample $F$ of $L_{m,n}$ has
$\|F\|_+>2^{m+n}$.
\end{enumerate}
Consequently, for $m=n$ one has $\|G_{n,n}\|=12n+9$, $\bar\theta_{G_{n,n}}=2^{n+1}-1$
and $\|F\|_+>\bar\theta_{G_{n,n}}^{\,2}/4$; in particular, for this typing at
fan-out one there is no bound of the form $p(\|G\|)\,\bar\theta_G$ with $p$ a
polynomial, in contrast with Theorem~\ref{thm:local-group-polydata}.
\end{proposition}

The proof is in Appendix~\ref{app:u1-typings}.  The pure word
$\mathtt a^{2^{m+n}}$ is the product of the block length and the number of
blocks, while $\bar\theta_G$ is about their sum, since a shortest sentence using
the long block fills every other position with $\mathtt b$.  Words containing a
letter other than $\mathtt a$ and $\mathtt b$ have empty distribution in
$L_{m,n}$, so the proposition holds verbatim for every typing onto $U_1$ in
which $\mathtt a$ is the only letter of type $1$ and $\mathtt b$ has type $\zeta$.
The data for these typings at fan-out one are nevertheless polynomial
(Corollary~\ref{cor:two-letter-fanout-one} below), and for $s_{U,2}$ so they are
at every fan-out $f\ge2$ (Corollary~\ref{cor:two-letter-semilattice}).

\paragraph{Isolated letters and two letters.}
The family of Theorem~\ref{thm:fanout-one-data}(iii) uses three letters, two of
them of type $1$.  A letter whose type is realized only by its own powers
already yields polynomial data, over any alphabet.

\begin{definition}[Isolated letter]
\label{def:isolated-letter}
A letter $\mathtt a\in\Sigma$ is \emph{isolated} for a typing $s$ if
$s^{-1}(s(\mathtt a))=\mathtt a^+$ and all words that contain $\mathtt a$ but do
not lie in $\mathtt a^+$ have one common type.
\end{definition}

For instance, let $\Sigma=\{\mathtt a,\mathtt b\}$ and let $c$ be the
\emph{content typing}, $c(w)$ the set of letters of $w$, onto the
three-element semilattice $\{\{\mathtt a\},\{\mathtt b\},\Sigma\}$ under union;
both letters are isolated for $c$, since a word that contains $\mathtt a$ and
lies outside $\mathtt a^+$ has content $\Sigma$.  If $s:\Sigma^+\twoheadrightarrow
U_1$ and $\mathtt a$ is the only letter of type $1$, then $s^{-1}(1)=\mathtt a^+$
and all other words have type $\zeta$, so $\mathtt a$ is isolated, over any
alphabet.  Over three or more letters no letter is isolated for the content
typing, since $\mathtt{ab}$ and $\mathtt{ac}$ have different contents.

\begin{theorem}[Isolated letters: short pure representatives at fan-out one]
\label{thm:two-letter-fanout-one}
Let $\mathtt a$ be isolated for a typing $s:\Sigma^+\twoheadrightarrow\mathsf S$,
let $L\in\Cmcf{1}{s}$, and let $G$ be a reduced good presentation of
$L\cap\Sigma^+$ of fan-out one and any rank.  For every nonterminal $A$ of $G$,
or of its binarization $G'$ \textup(Appendix~\ref{app:two-letter-fanout-one}\textup),
that has a yield in $\mathtt a^+$, the typed nonterminal $A_{s(\mathtt a)}$
satisfies
\[
 \tau_{A_{s(\mathtt a)}}\le C(\|G\|+2)^{22}\bar\theta_G^{\,8},
\]
where $C<10^{43}$ is an absolute constant, independent of $\Sigma$, $s$ and the
rank.  In particular, if $L$ contains a word of $\mathtt a^+$, its shortest one
has length at most $C(\|G\|+2)^{22}\bar\theta_G^{\,8}$.
\end{theorem}

The theorem bounds only the type of the isolated letter.  The other types are
bounded, for the two cases of the next corollary, by
Lemma~\ref{lem:app-h-other}, and the assumption that all words containing
$\mathtt a$ outside $\mathtt a^+$ have one type is essential: it fails for the
content typing over three or more letters, where the data can be exponential
(see below).

\begin{corollary}[Polynomial data from isolated letters]
\label{cor:two-letter-fanout-one}
Let $r_{\max}\ge2$ and let $s$ be
\begin{enumerate}[label=\textup{(\roman*)},leftmargin=*]
\item a typing over $\Sigma=\{\mathtt a,\mathtt b\}$ whose positive image is a
semilattice \textup(for instance $c$ or $s_{U,2}$\textup), or
\item a typing $s:\Sigma^+\twoheadrightarrow U_1$, over any finite alphabet,
in which exactly one letter has type $1$.
\end{enumerate}
Every $L\in\Cmcf{1}{s}$ with a reduced good fan-out-one presentation $G$ of
$L\cap\Sigma^+$, of any rank, has a common characteristic sample $K$ for
$\mathcal B^{\le r_{\max}}_{1,s}$ and $\mathcal B_{1,s}$ with
\[
 \|K\|_+\le C'r_{\max}(\|G\|+2)^{23}\bar\theta_G^{\,8}+\mathbf 1_{\{\eps\in L\}},
\]
where $C'<10^{45}$ is an absolute constant.
\end{corollary}

\begin{corollary}[Typings onto $U_1$ at fan-out one]
\label{cor:u1-classification}
Let $r_{\max}\ge2$ and let $s:\Sigma^+\twoheadrightarrow U_1$ be a typing.  Then
the bound required of each typing in Definition~\ref{def:uniform-polydata}
holds for $s$ at fan-out one and rank $r_{\max}$ if and only if exactly one
letter has type $1$, that is, $s^{-1}(1)=\mathtt a^+$ for some letter $\mathtt a$.  More precisely:
\begin{enumerate}[label=\textup{(\roman*)},leftmargin=*]
\item if exactly one letter has type $1$, the data are bounded as in
Corollary~\ref{cor:two-letter-fanout-one}, for presentations of every rank, and
by Proposition~\ref{prop:binary-u1-quadratic} not by any $p(\|G\|)\bar\theta_G$
with $p$ a polynomial;
\item if at least two letters have type $1$, the targets of
Corollary~\ref{cor:typing-level-obstruction} need data exponential in the size
and the thickness of their reduced good rank-two presentations.
\end{enumerate}
\end{corollary}

\begin{proof}
Since $s$ is onto and $U_1$ is generated by the types of the letters, some letter
has type $1$ and some letter has type $\zeta$.  (i) is
Corollary~\ref{cor:two-letter-fanout-one}(ii) together with
Proposition~\ref{prop:binary-u1-quadratic} and the remark following it, applied
with $\mathtt b$ a letter of type $\zeta$.  For (ii) let $\mathtt a\ne\mathtt b$ have
type $1$ and $\mathtt c$ type $\zeta$; the idempotents $e:=1$ and $z:=\zeta$ are
distinct with $ez=ze=z$, and $u:=\mathtt a$, $v:=\mathtt b$, $w:=\mathtt c$ are
letter-disjoint with $s(u)=s(v)=e$ and $s(w)=z$, so
Corollary~\ref{cor:typing-level-obstruction} applies.
\end{proof}

The proofs of Theorem~\ref{thm:two-letter-fanout-one} and
Corollary~\ref{cor:two-letter-fanout-one} are in
Appendix~\ref{app:two-letter-fanout-one}.  The idea is as follows.  Take a
node-minimal derivation $T$ of a shortest yield $\mathtt a^n$ of a nonterminal.
Two occurrences of one label at distance $d$ give, by cancellation of pure
blocks, a base $z_0$ outside $\mathtt a^+$ that commutes with $\mathtt a^d$ and
satisfies $z_0z_0\sim\mathtt a^{c_0}z_0$.  Short yields outside $\mathtt a^+$ of the
labels on the heavy path of $T$, padded by at most $d$ letters on each side,
collapse to $\mathtt a^ez_0$; isolation is used exactly here, since it makes all
words that contain $\mathtt a$ and lie outside $\mathtt a^+$ comparable.  The
set of exponents $f$ with $\mathtt a^fz_0$ in a given distribution class is a
residue class modulo some $p_0$; if $p_0$ exists, then $\mathtt a\sim\mathtt
a^{1+p_0}$, so pure words have short equivalent representatives.  Otherwise,
membership is decided by the values of the codes.  Each code is chosen to exceed
the sum of the previous ones, which costs a factor of at most $3$ per code, and
the permutations of these codes are counted against the size of the grammar,
as for the doubling families.

Corollary~\ref{cor:two-letter-fanout-one} does not contradict
Theorem~\ref{thm:fanout-one-data}(i), whose quantification is over all
alphabets: the semilattices form a pseudovariety not contained in
$\mathbf{LG}$, and the failure occurs over three letters.  Since the positive
image of every two-letter typing onto a semilattice is a quotient of the image
of $c$, part~(i) of the corollary covers all of them.  Over three or more
letters the failure of isolation for the content typing $c_\Sigma$ is
necessary: $L_n\in\Cmcf{1}{c_\Sigma}$ for $\Sigma\supseteq\{\mathtt a,\mathtt b,\mathtt
c\}$ by Proposition~\ref{prop:three-letter-family}(ii), and $W_n$ is its only word
of content $\{\mathtt a,\mathtt b\}$, so Proposition~\ref{prop:rare-type} forces data
of length at least $3\cdot2^n-1$.  In the proof the failure occurs when the codes
are collapsed: a short yield of content $\{\mathtt a,\mathtt c\}$ and a yield of
content $\{\mathtt a,\mathtt b\}$ of one nonterminal need not be equivalent.
Thus, at fan-out one, a typing onto $U_1$ has polynomial data if and only if
exactly one letter has type $1$ (Corollary~\ref{cor:u1-classification}); the data
are then polynomial but not of the form $p(\|G\|)\,\bar\theta_G$
(Proposition~\ref{prop:binary-u1-quadratic}).

Table~\ref{tab:scope} collects, for each main claim, the fixed parameters,
the operators and the side conditions.

% [PART 2] Semilattice reduction and two-letter results formerly in Section 8
\paragraph{Reduction to semilattices.}
Let $\mathbf{SL}$ be the pseudovariety of finite semilattices (commutative
idempotent semigroups).  For a nonempty finite alphabet $\Sigma$, the
\emph{content typing} $c_\Sigma$ maps $w\in\Sigma^+$ to the set of letters
occurring in $w$, a morphism onto the nonempty subsets of $\Sigma$ under union.

\begin{lemma}[Coarser typings inherit characteristic samples]
\label{lem:typing-monotone}
Let $s\preceq t$ be typings, $f,r\ge1$, and $L\in\Cmcf{f}{s}$.  For every
finite $K\subseteq L$,
\[
 L\bigl(\mathcal B^{\le r}_{f,t}(K)\bigr)\subseteq
 L\bigl(\mathcal B^{\le r}_{f,s}(K)\bigr)\subseteq L,
\]
and likewise for $\mathcal B_{f,t}$ and $\mathcal B_{f,s}$.  Hence every
characteristic sample of $L$ for $\mathcal B^{\le r}_{f,t}$ (resp.\
$\mathcal B_{f,t}$) is one for $\mathcal B^{\le r}_{f,s}$ (resp.\
$\mathcal B_{f,s}$).
\end{lemma}

\begin{proof}
Both grammars have the same nonterminals and the same Start, Const and Comp
rules (Definition~\ref{def:reconstruction}).  A Link rule of the $t$-grammar
joins observed tuples with a common context and equal $t$-types; since
$\ker t\subseteq\ker s$, their $s$-types are equal, so it is a Link rule of the
$s$-grammar.  This gives the first inclusion; the second is
Proposition~\ref{prop:reconstruction-soundness}.  If $F$ is characteristic for
the $t$-operator and $F\subseteq K\subseteq L$, the left-hand side equals $L$,
hence so does the middle term.
\end{proof}

\begin{proposition}[Reduction of the fan-out boundary to semilattices]
\label{prop:sl-reduction}
Fix $f\ge1$ and $r_{\max}\ge1$, and read uniform polynomial characteristic-data
length at fan-out $f$ and rank $r_{\max}$
(Definition~\ref{def:uniform-polydata}).
\begin{enumerate}[label=\textup{(\roman*)},leftmargin=*]
\item Every pseudovariety $\mathbf V\not\subseteq\mathbf{LG}$ contains
$\mathbf{SL}$.  Consequently, the pseudovarieties with uniform polynomial
characteristic-data length are exactly those contained in $\mathbf{LG}$ if and
only if $\mathbf{SL}$ does not have it.
\item $\mathbf{SL}$ has uniform polynomial characteristic-data length if and
only if, for every nonempty finite alphabet $\Sigma$, the content typing
$c_\Sigma$ satisfies the condition of Definition~\ref{def:uniform-polydata}.
\end{enumerate}
\end{proposition}

\begin{proof}
(i) A finite semigroup $\mathsf S\notin\mathbf{LG}$ is the positive image of
its evaluation typing over the alphabet $\mathsf S$, so it has a subsemigroup
isomorphic to $U_1$ (Theorem~\ref{thm:lg-characterization}); hence
$U_1\in\mathbf V$.  Every finite semilattice $\mathsf H$ embeds into
$U_1^{\mathsf H}$ via $x\mapsto(\chi_z(x))_{z\in\mathsf H}$, where $\chi_z(x)=1$ if
$zx=z$ and $\chi_z(x)=\zeta$ otherwise (a morphism, since $zxy=z$ iff $zx=z$ and
$zy=z$), so $\mathbf{SL}\subseteq\mathbf V$.  The property of
Definition~\ref{def:uniform-polydata} passes from a pseudovariety to its
subpseudovarieties, because fewer typings are quantified.  If $\mathbf{SL}$
lacks it, so does every $\mathbf V\not\subseteq\mathbf{LG}$, while every
$\mathbf V\subseteq\mathbf{LG}$ has it by Theorem~\ref{thm:local-group-polydata}.
Conversely, $U_1\notin\mathbf{LG}$, so $\mathbf{SL}\not\subseteq\mathbf{LG}$.

(ii) The image of $c_\Sigma$ is a semilattice, which gives one direction.  For
the other, let $s:\Sigma^+\twoheadrightarrow\mathsf S$ with $\mathsf S\in\mathbf{SL}$.
Since $\mathsf S$ is commutative and idempotent, $s(w)$ is the product of
$s(a)$ over the letters $a$ occurring in $w$, so $s\preceq c_\Sigma$ and
$\Cmcf{f}{s}\subseteq\Cmcf{f}{c_\Sigma}$ (item~(a) after
Definition~\ref{def:class-cmcf}).  For $L\in\Cmcf{f}{s}$ and a presentation $G$,
a characteristic sample for $\mathcal B^{\le r_{\max}}_{f,c_\Sigma}$ of length at
most $p_{c_\Sigma}(\|G\|,\bar\theta_G)$ is one for
$\mathcal B^{\le r_{\max}}_{f,s}$ by Lemma~\ref{lem:typing-monotone}; take
$p_{s}:=p_{c_\Sigma}$.
\end{proof}

By Theorem~\ref{thm:fanout-one-data}, $\mathbf{SL}$ lacks the property at
fan-out one when $r_{\max}\ge2$, although over two letters every typing in
$\mathbf{SL}$ has polynomial data there (Corollary~\ref{cor:two-letter-fanout-one}).
At fan-out two the question is open (Open Problem~\ref{op:fanout-two}), and the
next theorem settles two-letter alphabets at every fan-out $f\ge2$.  Rare-type lower bounds transfer to content
typings: for $s\preceq c_\Sigma$ every content type contained in an $s$-type $m$
has shortest word at least $\mu_m(L)$, computed for $s$.

\paragraph{Two letters.}
Call $L$ \emph{Parikh-saturated} if membership of a word of $\Sigma^+$ in $L$
depends only on its Parikh vector.

\begin{theorem}[Content typings on two letters at every fan-out $f\ge2$]
\label{thm:two-letter-content}
Let $\Sigma=\{\mathtt a,\mathtt b\}$, $c:=c_\Sigma$, $f\ge2$, $r_{\max}\ge2$, and put
$D_f:=(6f+1)3^f$ and
$C_{f,r_{\max}}:=(r_{\max}+1)\bigl(4r_{\max}(3f+1)+D_f\bigr)+3$, and let
$H_m:=\sum_{k=1}^m1/k\le1+\ln m$ \textup($H_0:=0$\textup).  For every
$L\in\Cmcf{f}{c}$ and every reduced good fan-out-$\le f$, rank-$\le r_{\max}$
presentation $G$ of $L\cap\Sigma^+$,
$\tau(G,c)\le C_{f,r_{\max}}|V|^2\bar\theta_G^{\,4}H_{\bar\theta_G}$, and the operators
$\mathcal B^{\le r_{\max}}_{f,c}$ and $\mathcal B_{f,c}$ have a common
characteristic sample $K$ with
\[
 \|K\|_+\le4^{fr_{\max}}\bigl(3+r_{\max}C_{f,r_{\max}}\bigr)\|G\|^3\bar\theta_G^{\,4}H_{\bar\theta_G}
 +\mathbf 1_{\{\eps\in L\}}.
\]
If $L$ is not Parikh-saturated, then
$t_{\mathrm{typ}}(G,c)\le D_f(r_{\max}+1)|V|\bar\theta_G$ and one may take
$\|K\|_+\le4^{fr_{\max}}\bigl(3+r_{\max}D_f(r_{\max}+1)\bigr)\|G\|^2\bar\theta_G
+\mathbf 1_{\{\eps\in L\}}$.
\end{theorem}

\begin{corollary}[Semilattice typings on two letters]
\label{cor:two-letter-semilattice}
Let $f\ge2$, $r_{\max}\ge2$ and let $s:\{\mathtt a,\mathtt b\}^+\twoheadrightarrow\mathsf S$ be a
typing with $\mathsf S\in\mathbf{SL}$, for instance $s_{U,2}$
\textup($\mathtt a\mapsto1$, $\mathtt b\mapsto\zeta$\textup).  Then the bound
required of each typing in Definition~\ref{def:uniform-polydata} holds for this
fixed typing at fan-out $f$ and rank $r_{\max}$, with the bound of
Theorem~\ref{thm:two-letter-content}, for
$\mathcal B^{\le r_{\max}}_{f,s}$ and for $\mathcal B_{f,s}$.
\end{corollary}

\begin{proof}
As in the proof of Proposition~\ref{prop:sl-reduction}(ii), $s\preceq c$ and
$\Cmcf{f}{s}\subseteq\Cmcf{f}{c}$; apply Theorem~\ref{thm:two-letter-content} and
Lemma~\ref{lem:typing-monotone}.
\end{proof}

The proof is in Appendix~\ref{app:two-letter-content}.  It uses the
conditions of arities one and two, which hold at every fan-out $f\ge2$, through
runs of one letter: once $x^4$ occurs in a word of
$L$ for a letter $x$, membership depends only on the $x$-skeleton (each maximal
run of $x$ replaced by one $x$) and the number of letters $x$, and two runs of the other letter next to a run $xx$ either
force Parikh saturation or merge into a forbidden longer run.  Without
saturation this bounds the antichains of heavy nodes in a shortest typed
derivation; with saturation, node-minimal derivations of words with a
prescribed Parikh vector are analysed directly.  The theorem does not separate
fan-outs: by Corollary~\ref{cor:two-letter-fanout-one}, two-letter
semilattice typings have polynomial data at fan-out one as well, so over two
letters they have polynomial data at every fan-out.  For three or more
letters a long component of content $\{\mathtt a,\mathtt b\}$ need not contain a
long run, and the question is open.

% [PART 3] Full Appendix H/I/J proofs (isolated letters, U1 typings, two-letter content)
\section{Isolated Letters at Fan-Out One}
\label{app:two-letter-fanout-one}

This appendix proves Theorem~\ref{thm:two-letter-fanout-one} and
Corollary~\ref{cor:two-letter-fanout-one}.  Throughout,
$s:\Sigma^+\twoheadrightarrow\mathsf S$ is a typing, $\mathtt a\in\Sigma$ is isolated
for $s$ (Definition~\ref{def:isolated-letter}), and $L\in\Cmcf{1}{s}$; we may
assume $\eps\notin L$.  A word of $\Sigma^+$ is \emph{pure} if it lies in
$\mathtt a^+$, \emph{$\mathtt a$-free} if it does not contain $\mathtt a$, and
\emph{$\mathtt a$-mixed} otherwise, that is, if it contains $\mathtt a$ and another
letter.  Isolation is used only through two facts:
\begin{enumerate}[label=\textup{(I\arabic*)},leftmargin=*]
\item all pure words have the type $s(\mathtt a)$, and the words of this type are
the pure words;
\item all $\mathtt a$-mixed words have one type.
\end{enumerate}
For $x\in\Sigma^+$ write $\D(x):=\D^{(1)}_L(x)$ and $x\sim y$ if
$\D(x)=\D(y)$; this is a congruence on $\Sigma^+$, and words with empty
distribution are $\sim$-equivalent.  Thus $x\equiv_{L,s}y$ if and only if
$x\sim y$ and $s(x)=s(y)$.

\subsection{Preliminaries}

\begin{lemma}[Cancellation]
\label{lem:app-h-cancel}
Let $x,y\in\Sigma^+$ have the same type, $z\in\Sigma^+$, $\D(zx)\ne\emptyset$
and $zx\sim zy$.  Then $x\sim y$.  The same holds on the right.
\end{lemma}

\begin{proof}
For $(l,r)\in\D(zx)$ we get $lzyr\in L$, so $x$ and $y$ share the context
$(lz,r)$.
\end{proof}

We apply the lemma only to pure $x,y$ or to $\mathtt a$-mixed $x,y$, which have
the same type by (I1), respectively (I2).  If a nonterminal $X$ has a
derivational context $(l,r)$, two yields of $X$ of the same type share $(l,r)$
and are $\sim$-equivalent (Lemma~\ref{lem:typed-homogeneity}).  For a type $t$
realized by $X$ we write $\tau(X,t)$ for the representative length of the typed
nonterminal $X_t$, the least length of a word of type $t$ that is
$\sim$-equivalent to the yields of $X$ of type $t$.

\paragraph{Binarization.}
Let $G$ be a reduced good presentation of $L$ of fan-out one, and let
$\sigma_G$ be the number of letters occurring in its rules, so
$\sigma_G\le\min\{|\Sigma|,\|G\|\}$.  Replace every rule $X\to Y_1\cdots Y_t$
with $t\ge3$ by a right-linear chain through fresh nonterminals, replace
terminals occurring in right sides of length at least two by preterminals
$T_{\mathtt x}\to\mathtt x$ (only for the letters $\mathtt x$ that are used), and trim.
The resulting grammar $G'$ is reduced and good, has rules of the forms
$X\to YZ$, $X\to Y$ and $X\to\mathtt x$ with $\mathtt x\in\Sigma$, and generates $L$.
A rule $X\to\gamma$ of $G$ gives at most $\max\{0,|\gamma|-2\}$ fresh
nonterminals and at most $\max\{1,|\gamma|-1\}$ rules, both at most the size
$1+r+|\gamma|$ of the rule in $\|G\|$, so
\[
 k_V:=|V'|\le\|G\|+\sigma_G\le2\|G\|,\qquad |P'|\le\|G\|+\sigma_G,
\]
and $k_V\le\|G\|+2$, $|P'|\le\|G\|+2$ if $|\Sigma|=2$.  Moreover
$\theta:=\bar\theta_{G'}\le\bar\theta_G$, since a derivation using a new rule
corresponds to one using the rule of $G$ it comes from, and a derivation using a
preterminal rule $T_{\mathtt x}\to\mathtt x$ to one using a rule of $G$ that writes
$\mathtt x$.  Every nonterminal of
$G'$ has a yield of length $\le\theta$ and a derivational context $(l,r)$ with
$|l|+|r|\le\theta$: take a shortest sentence through a rule with that
nonterminal on the left, respectively on the right.  The nonterminals of $G$
are nonterminals of $G'$ with the same yields.

\subsection{The Tree and the Base}

Fix a nonterminal $A$ of $G'$ with a pure yield; let $\mathtt a^n$ be the
shortest one, $T$ a derivation tree of $\mathtt a^n$ from $A$ with the least
number of nodes, and $(l,r)$ a derivational context of $A$ with
$|l|+|r|\le\theta$.  Replacing subtrees of $T$ by derivations from the same
labels gives yields $w$ of $A$, hence $lwr\in L$.  A node is \emph{heavy} if
its yield is longer than $\theta$, and a heavy node without heavy children is a
\emph{seed}.

\begin{lemma}
\label{lem:app-h-tree}
\begin{enumerate}[label=\textup{(\alph*)},leftmargin=*]
\item No label repeats on a root-to-leaf path of $T$; such a path has at
most $k_V$ nonterminal nodes.
\item All nodes of $T$ with the same label $Y$ have the same yield
$\mathtt a^{n_Y}$.
\item Every label $Y$ of a heavy node has a yield $\omega_Y\notin\mathtt a^+$ of
length $\le\theta$.
\item Seeds have yields of length $\le2\theta$; seeds are pairwise disjoint;
a subtree containing $k\ge1$ seeds has a yield of length $\le k(k_V+2)\theta$; and
the left ends of two consecutive seeds differ by at most $2k_V\theta$.
\end{enumerate}
\end{lemma}

\begin{proof}
(a) If an ancestor and a descendant share a label, replacing the ancestor's
subtree by the descendant's gives a shorter pure yield of $A$ or a smaller
tree.  (b) Otherwise replace the longer subtree by the shorter one.  (c) Let
$\rho$ be the rule at a heavy node labelled $Y$.  A sentence of length
$\le\theta$ using $\rho$ contains a yield $y$ of $Y$ with $|y|\le\theta<n_Y$;
if $y\in\mathtt a^+$, inserting it into $T$ gives a shorter pure yield of $A$.
(d) A seed is not a unit node and has two light children.  If a seed $v$ lay
below a seed $u$, the child of $u$ above $v$ would be heavy.  The heavy nodes of
a subtree with $k$ seeds lie on paths to seeds, so there are at most $kk_V$ of
them, each non-seed one with at most one light child.  Between consecutive
seeds $u,v$ lie the yields of the subtrees hanging off the two paths from
their least common ancestor, at most $2(k_V-1)$ of them; each is light, since a
heavy one would contain a seed between $u$ and $v$.
\end{proof}

Let $X_0$ be the root of $T$ and $X_{i+1}$ the child of $X_i$ with the longer
yield, so that $n_{X_i}\le2n_{X_{i+1}}$.  If the root subtree has no seed, the
root is light and $n\le\theta$; if it has at least one and at most $k_V$ seeds,
then $n\le k_V(k_V+2)\theta$.  Otherwise let $i^*$ be the largest index
whose subtree contains at least $k_V+1$ seeds; then $X_{i^*+1}$ exists,
$i^*\le k_V-2$ by Lemma~\ref{lem:app-h-tree}(a), and
$B_0:=n_{X_{i^*}}\le2k_V(k_V+2)\theta$; we may assume $n_{X_{i^*}}<n$, and
$k_V\ge2$ since $T$ has a binary node.  The seeds below $X_{i^*}$ are
consecutive in $T$; among $k_V+1$ consecutive ones two have the same label $Y$,
at positions $x_1<x_2$ with $d:=x_2-x_1\le2k_V^2\theta$.  Put $g_0:=d-n_Y\ge0$,
$\omega:=\omega_Y$, $x:=x_1$, $y:=n-x_2-n_Y$, and $y_1:=n-x_1-n_Y=d+y$.
Replacing both seeds, the first, or the second by $\omega$ gives the yields
\[
 W_1=\mathtt a^x\omega\mathtt a^{g_0}\omega\mathtt a^y,\qquad
 W_2=\mathtt a^x\omega\mathtt a^{d+y},\qquad
 W_3=\mathtt a^{x+d}\omega\mathtt a^y
\]
of $A$.  Here $x+g_0+y\ge1$: otherwise the two seeds would cover the yield of
the root, which contains at least $k_V+1\ge3$ seeds.  Since $d\ge n_Y\ge1$ and
$\omega\notin\mathtt a^+$, the words $W_1,W_2,W_3$ are $\mathtt a$-mixed.

\begin{lemma}[Base]
\label{lem:app-h-base}
Put $z_0:=\omega\mathtt a^{g_0}$ and $c_0:=d$ if $g_0\ge1$ or $\omega$ is
$\mathtt a$-mixed, and $z_0:=\omega\mathtt a^d$ and $c_0:=2d$ if $g_0=0$ and $\omega$
is $\mathtt a$-free.  Then $z_0$ is $\mathtt a$-mixed, $\D(z_0)\ne\emptyset$,
$d\mid c_0$,
\[
 \mathtt a^dz_0\sim z_0\mathtt a^d\qquad\text{and}\qquad z_0z_0\sim\mathtt a^{c_0}z_0.
\]
\end{lemma}

\begin{proof}
By (I2), $W_1,W_2,W_3$ have one type; they share $(l,r)$, hence are
$\sim$-equivalent.  Cancelling $\mathtt a^x$ on the left and $\mathtt a^y$ on the
right of $W_2\sim W_3$ (Lemma~\ref{lem:app-h-cancel}; the remaining words contain
$\mathtt a^d$ and $\omega$, so they are $\mathtt a$-mixed, and the words cancelled
from are factors of $W_2$, so their distributions are nonempty) gives
$\omega\mathtt a^d\sim\mathtt a^d\omega$.  If $g_0\ge1$
or $\omega$ is $\mathtt a$-mixed, the same cancellation in $W_1\sim W_2$ (now
$\omega\mathtt a^{g_0}\omega$ is $\mathtt a$-mixed) gives $\omega\mathtt
a^{g_0}\omega\sim\omega\mathtt a^d$, whence $z_0z_0\sim\omega\mathtt
a^{d+g_0}=z_0\mathtt a^d$ and $\mathtt a^dz_0\sim z_0\mathtt a^d$.  If $g_0=0$ and
$\omega$ is $\mathtt a$-free, then $n_Y=d$ and $x+y\ge1$.  For $y\ge1$,
cancelling $\mathtt a^x$ and $\mathtt a^{y-1}$ gives $\omega\omega\mathtt
a\sim\omega\mathtt a^{d+1}$, so $z_0z_0\sim\omega\omega\mathtt a^{2d}\sim\omega\mathtt
a^{3d}=z_0\mathtt a^{2d}$; for $y=0$, cancelling $\mathtt a^{x-1}$ gives $\mathtt
a\omega\omega\sim\mathtt a\omega\mathtt a^d$, so, keeping one letter $\mathtt a$
in front of $\omega\omega$, $z_0z_0\sim\mathtt a^{2d}\omega\omega=\mathtt
a^{2d-1}(\mathtt a\omega\omega)\sim\mathtt a^{2d-1}(\mathtt a\omega\mathtt
a^d)=\mathtt a^{2d}z_0$; the $\mathtt a$-free word $\omega\omega$ is never compared
with anything.  The words compared here contain
$\mathtt a$ and $\omega$, so they are $\mathtt a$-mixed, and those cancelled from
are factors of $W_1$.  In both cases $z_0\mathtt
a^{2d}\sim\mathtt a^{2d}z_0$ and $\mathtt a^dz_0\sim z_0\mathtt a^d$.  Finally $z_0$
is a factor of $W_1$ or of $W_2$.
\end{proof}

For a typing onto $U_1$ in which $\mathtt a$ is the only letter of type $1$, the
$\mathtt a$-free words have the type of the $\mathtt a$-mixed ones, and the
second case could be avoided; we do not use this.

An \emph{annotated code} is a pair $(z,e)$ with $e\ge0$, $d\mid e$ and $z\sim
\mathtt a^ez_0$; its \emph{value} is $e+c_0$.  A pure block $\mathtt a^R$ with
$d\mid R$ has value $R$.  For a specified product $u$ of annotated codes and
such blocks, $\operatorname{val}(u)$ is the sum of the values of its factors.
(The exponent $e$ is not determined by $z$ in general, so values are attached
to annotations.)

\begin{lemma}
\label{lem:app-h-collapse}
If $u$ contains at least one code, then $u\sim\mathtt
a^{\operatorname{val}(u)-c_0}z_0$.
\end{lemma}

\begin{proof}
Replace each code by $\mathtt a^ez_0$, move all blocks to the left with
$\mathtt a^dz_0\sim z_0\mathtt a^d$, and use $z_0^t\sim\mathtt a^{(t-1)c_0}z_0$,
which follows by induction from $z_0z_0\sim\mathtt a^{c_0}z_0$ and $d\mid c_0$.
\end{proof}

\subsection{Codes along the Heavy Path}

Let $i_0$ be the least index with $n_{X_i}<n$; then $1\le i_0\le i^*$.  For
$i_0\le i\le i^*$ let $(\mathtt a^{p_i},\mathtt a^{q_i})$ be the context of
$X_i$ inside $T$, so $p_i+q_i\ge1$; the subtree of $X_i$ contains both seeds;
let $L_i,R_i$ be the numbers of letters to the left and right of the first seed
inside it, so $R_i\ge g_0+n_Y\ge1$, and put $\Omega_i:=\mathtt a^{L_i}\omega\mathtt
a^{R_i}$, an $\mathtt a$-mixed yield of $X_i$.

\begin{lemma}
\label{lem:app-h-borrow}
If $\omega_{X_i}$ is $\mathtt a$-mixed, then $\omega_{X_i}\sim\Omega_i$.  If
$\omega_{X_i}$ is $\mathtt a$-free and $p_i\ge1$, then $\mathtt a\omega_{X_i}\sim\mathtt
a\Omega_i$; if $\omega_{X_i}$ is $\mathtt a$-free and $p_i=0$, then
$\omega_{X_i}\mathtt a\sim\Omega_i\mathtt a$.
\end{lemma}

\begin{proof}
Both $l\mathtt a^{p_i}\omega_{X_i}\mathtt a^{q_i}r$ and $l\mathtt a^{p_i}\Omega_i
\mathtt a^{q_i}r$ lie in $L$.  In the first case $\omega_{X_i}$ and $\Omega_i$ are
$\mathtt a$-mixed and share the context $(l\mathtt a^{p_i},\mathtt a^{q_i}r)$.  In the
second case $\mathtt a\omega_{X_i}$ and $\mathtt a\Omega_i$ are $\mathtt a$-mixed and
share the context $(l\mathtt a^{p_i-1},\mathtt a^{q_i}r)$; the third case is
symmetric, with $q_i\ge1$.  Conclude by (I2).
\end{proof}

Choose $s_i',s_i\in[0,d]$ with $d\mid L_i+s'_i$, with $R_i+s_i\equiv g_0\pmod d$
in the first case of Lemma~\ref{lem:app-h-base} and $R_i+s_i\equiv0\pmod d$ in
the second, and with $s'_i\ge1$, respectively $s_i\ge1$, when a letter is
borrowed in Lemma~\ref{lem:app-h-borrow}.  Put $\hat z_i:=\mathtt
a^{s'_i}\omega_{X_i}\mathtt a^{s_i}$.  By Lemma~\ref{lem:app-h-borrow},
$\hat z_i\sim\mathtt a^{s'_i+L_i}\omega\mathtt a^{R_i+s_i}$.  In the first case
$R_i+s_i=g_0+kd$ with $k\ge0$, and in the second $R_i+s_i=d+kd$ with $k\ge0$,
since $R_i\ge g_0+n_Y$; so $\hat z_i\sim\mathtt a^{e_i}z_0$ with
$e_i:=s'_i+L_i+kd$, and $(\hat z_i,e_i)$ is an annotated code, of value
\[
 v_i=n_{X_i}+s_i+s'_i,\qquad|\hat z_i|\le\theta+2d,
\]
in both cases (using $n_{X_i}=L_i+n_Y+R_i$ and $d-g_0=n_Y$, respectively
$n_Y=d$).  For $i<i_0$ put $v_i:=n$ (no code).  Since $X_i$ is an ancestor of
$X_{i^*}$, whose yield contains both seeds, $n_{X_i}\ge d+n_Y$ for $i\le i^*$.
Hence, for all $i\le i^*$,
\begin{equation}
\label{eq:app-h-star}
 v_i\ge n_{X_i}>d,\qquad v_i\le2v_{i+1}+2d\ \ (i<i^*),
\end{equation}
and $v_{i^*}\le B_0+2d$; moreover $n\le2n_{X_{i_0}}\le2v_{i_0}$.

The yield $\mathtt a^{x_1}\omega\mathtt a^{y_1}$ of $A$ equals $\mathtt
a^{x_1}z_0\mathtt a^{h}$ with $h:=y_1-g_0$, respectively $h:=y_1-d$.  Writing
$h=qd+\rho$ and $x_1+qd=e^*+\rho'$ with $0\le\rho,\rho'<d$ and $d\mid e^*$, we
obtain $l_\circ\mathtt a^{e^*}z_0r_\circ\in L$ for $l_\circ:=l\mathtt a^{\rho'}$
and $r_\circ:=\mathtt a^\rho r$, and $V^*:=e^*+c_0=n-\rho-\rho'$.

\subsection{Periodicity and Membership}

\begin{lemma}
\label{lem:app-h-period}
Let $P_0:=\{p\ge0:\mathtt a^pz_0\sim z_0\}$.
\begin{enumerate}[label=\textup{(\roman*)},leftmargin=*]
\item $P_0=\{0\}$ or $P_0=p_0\mathbb Z_{\ge0}$ for some $p_0\ge1$.
\item If $l'\mathtt a^fz_0r'\in L$ for some $f$, then
$\{f\ge0:l'\mathtt a^fz_0r'\in L\}$ is a single point if $P_0=\{0\}$ and a
full residue class modulo $p_0$ otherwise.
\item If $P_0\ne\{0\}$, then $\mathtt a\sim\mathtt a^{1+p_0}$, so every
$\mathtt a^k$ is $\equiv_{L,s}$-equivalent to some $\mathtt a^{k'}$ with $1\le k'\le p_0$.
\end{enumerate}
\end{lemma}

\begin{proof}
All words $\mathtt a^fz_0$ are $\mathtt a$-mixed, so they have one type by (I2).
(i) $P_0$ is closed under sums.  It is closed under differences $q-p$ with
$q>p>0$ in $P_0$ by Lemma~\ref{lem:app-h-cancel}, applied to
$\mathtt a^p(\mathtt a^{q-p}z_0)\sim\mathtt a^pz_0$, since $\D(\mathtt
a^pz_0)=\D(z_0)\ne\emptyset$; the case $p=0$ is trivial.
(ii) Two members $f<f'$ give $\mathtt a^fz_0\sim\mathtt a^{f'}z_0$, hence
$f'-f\in P_0$ by Lemma~\ref{lem:app-h-cancel}; conversely the set is closed
under adding and, by $\mathtt a^{p_0}z_0\sim z_0$, under subtracting $p_0$.  (iii)
For $(l',r')\in\D(z_0)$ both $l'\mathtt a^{p_0}z_0r'$ and $l'\mathtt
a^{2p_0}z_0r'$ lie in $L$, so the pure words $\mathtt a$ and $\mathtt a^{1+p_0}$,
which have one type by (I1), share the context $(l'\mathtt a^{p_0-1},z_0r')$.
\end{proof}

Let $D:=\operatorname{lcm}(d,p_0)$, and $D:=\infty$ if $P_0=\{0\}$.  By
Lemma~\ref{lem:app-h-collapse} and Lemma~\ref{lem:app-h-period}(ii), for
every specified product $u$ as above that contains a code,
\begin{equation}
\label{eq:app-h-member}
 l_\circ ur_\circ\in L\iff\operatorname{val}(u)\equiv V^*\pmod D,
\end{equation}
with equality when $D=\infty$; values and $V^*$ are multiples of $d$.  Two
annotations of the same word differ by a multiple of $D$, so membership does
not depend on the annotation.

\subsection{Counting}

\begin{lemma}
\label{lem:app-h-count}
Let $(z_1,e_1),\ldots,(z_M,e_M)$ be annotated codes, $M\ge6$, with
$|z_j|\le\ell_z$, values $w_j$ whose subset sums are pairwise distinct, and
$\Sigma_w:=\sum_jw_j<D$, and let $Q$ be a fixed specified product of codes and
blocks with $\Sigma_w+\operatorname{val}(Q)=V^*$.  Put $m:=|l_\circ|+|r_\circ|
+|Q|+\sum_j|z_j|$.  Then
\[
 (27/4)^{\lfloor M/3\rfloor}\le(M+1)^2k_V(m+1)^2(M\ell_z+1)^2.
\]
\end{lemma}

\begin{proof}
For a permutation $\pi$, the word $w_\pi:=l_\circ z_{\pi(1)}\cdots
z_{\pi(M)}Qr_\circ$ lies in $L$ by \eqref{eq:app-h-member}.  Take a
derivation tree of $w_\pi$ in $G'$, contract unit nodes, and give each node the
weight ``number of the codes $z_{\pi(1)},\ldots,z_{\pi(M)}$ whose first letter
lies in its yield''.  Leaves have weight $\le1$ and the root has weight $M$;
descending to the heavier child we reach a node $u$ of weight
$\kappa\in(M/3,2M/3]$, with label $X$ and yield interval $[p,q)$.  Let $I(\pi)$
be the set of indices $j$ such that $z_j$ lies inside $[p,q)$; then
$\kappa-1\le|I(\pi)|\le\kappa$.  The \emph{record} of $\pi$ consists of $X$,
$p$, $q$, and, for each end of $[p,q)$ lying strictly inside one of the $M$
codes, the index and the offset of that code; there are at most
$k_V(m+1)^2(M\ell_z+1)^2$ records.  The three blocks $l_\circ$, $z_{\pi(1)}\cdots
z_{\pi(M)}$ and $Qr_\circ$ occupy the same positions for all $\pi$.  The codes
whose first letter lies in $[p,q)$ are those of $I(\pi)$ and at most one code
that starts inside and ends after $q$; a code through $p$ starts before $p$.
This gives $\kappa-1\le|I(\pi)|$.  Each end of $[p,q)$ lies in $l_\circ$, in
$Qr_\circ$, on a boundary between two codes, or strictly inside a code; only in
the last case is a code cut there, and then its index and offset are part of the
record.  So the record also determines which of these four cases occurs at each
end.  No code is cut at both ends, since then $[p,q)$ would lie inside one code
and $\kappa\le1$.  The codes are pairwise distinct words: if $z_j=z_k$, their
annotations differ by a multiple of $D$, so $w_j\equiv w_k\pmod D$, hence
$w_j=w_k$ since both lie in $[0,D)$, and $j=k$ since the subset sums are
distinct.  (The splicing below works with positions and does not need this.)

Let $\pi,\pi'$ have the same record.  Replacing the subtree of $u$ in the tree
of $w_\pi$ by the corresponding subtree of $w_{\pi'}$ gives a derivation of
$w_\pi[0,p)\,w_{\pi'}[p,q)\,w_\pi[q,|w_\pi|)\in L$.  Since an end inside a code
has the same index and offset on both sides, the fragments join into complete
copies of the same code, and the word is $l_\circ u''Qr_\circ$ where $u''$ is a
product of complete codes in which $z_j$ occurs
$1-\mathbf 1[j\in I(\pi)]+\mathbf 1[j\in I(\pi')]$ times; this also covers ends
in $l_\circ$ or $Qr_\circ$, and both ends cannot lie in $Qr_\circ$ because
$\kappa\ge2$.  By \eqref{eq:app-h-member},
$\sum_{I(\pi')}w_j\equiv\sum_{I(\pi)}w_j\pmod D$; both sums lie in
$[0,\Sigma_w)\subseteq[0,D)$, so they are equal, and $I(\pi')=I(\pi)$.

Hence a record and $I=I(\pi)$, $|I|=j$, determine $\pi$ up to the order inside
($\le j!$) and the arrangement of the remaining codes around this block
($\le(M-j+1)!$).  Since $j\in\{\kappa-1,\kappa\}$ and $M/3<\kappa\le2M/3$, $j$
lies in $J:=[\lfloor M/3\rfloor,\lfloor2M/3\rfloor]$, and
$\binom{M+1}{j}\le(M+1)k_V(m+1)^2(M\ell_z+1)^2$ for some $j\in J$.  Put
$k:=\lfloor M/3\rfloor$.  As $\binom{M+1}{\cdot}$ is symmetric and unimodal and
$\lfloor M/3\rfloor+\lfloor2M/3\rfloor\le M$, every $j\in J$ satisfies
$\binom{M+1}{j}\ge\binom{M+1}{k}\ge\binom{3k}{k}$.  Finally, the numbers
$t_i:=\binom{3k}{i}3^{-i}(2/3)^{3k-i}$, $0\le i\le3k$, sum to $1$, and
$t_{i+1}/t_i=(3k-i)/(2i+2)$ is at least $1$ if and only if $i\le k-1$; so $t_k$
is the largest of them and $t_k\ge1/(3k+1)\ge1/(M+1)$.  Since
$t_k=\binom{3k}{k}4^k/27^k$, this gives $\binom{3k}{k}\ge(27/4)^k/(M+1)$, and the
lemma follows.
\end{proof}

\subsection{Selection and Filler}

Put $\Lambda:=\min(V^*,D)$, so $\Lambda\le V^*\le n$.

\begin{lemma}[Selection]
\label{lem:app-h-select}
There are $M\ge0$ indices $i^*=j_1>j_2>\cdots>j_M\ge i_0$ whose values
$w_t:=v_{j_t}$ are superincreasing \textup($w_t>w_1+\cdots+w_{t-1}$\textup), with
$w_1+\cdots+w_M<\Lambda$ and $\Lambda\le3^M(v_{i^*}+d)$.
\end{lemma}

\begin{proof}
If $v_{i^*}\ge\Lambda$, take $M:=0$.  Otherwise take $j_1:=i^*$.  Given
$j_t$ and $S_t:=w_1+\cdots+w_t<\Lambda$, let $i$ be the largest index below $j_t$
with $v_i>S_t$; it exists, since $v_0=n\ge\Lambda$ (recall $i_0\ge1$).  If
$S_t+v_i\ge\Lambda$, stop with $M:=t$; otherwise put $j_{t+1}:=i$.  The values
are superincreasing and their sum is below $\Lambda$ by construction.  A taken
value is less than $\Lambda\le n$, whereas $v_i=n$ for $i<i_0$; so all taken
indices are at least $i_0$.  By the choice of $i$, $v_{i+1}\le S_t$: for
$i+1<j_t$ by maximality, and for $i+1=j_t$ because $v_{j_t}=w_t$.  By
\eqref{eq:app-h-star}, $v_i\le2S_t+2d$.  Hence $S_{t+1}+d\le3(S_t+d)$ whenever $i$
is taken, and $\Lambda\le S_M+v_i<3(S_M+d)$ when the process stops.  With
$S_1=v_{i^*}$ this gives $\Lambda<3^M(v_{i^*}+d)$.
\end{proof}

\begin{lemma}[Filler]
\label{lem:app-h-filler}
Let $0<X<n$ be a multiple of $d$.  There is a specified product $Q$ of at most
$3k_V$ codes among $\hat z_{i_0},\ldots,\hat z_{i^*}$ and at most one pure block
$\mathtt a^{X'}$ with $d\mid X'$ and $X'<v_{i^*}$, such that
$\operatorname{val}(Q)=X$ and $|Q|\le(3k_V+2)\ell_z$, where
$\ell_z:=5k_V^2\theta$.
\end{lemma}

\begin{proof}
For $i=i_0,\ldots,i^*$ in turn, append $q_i:=\lfloor X/v_i\rfloor$ copies of
$\hat z_i$ and replace $X$ by $X-q_iv_i$.  Then $q_{i_0}\le1$, since
$X<n\le2v_{i_0}$.  For $i>i_0$ the current $X$ is less than $v_{i-1}$, and
$v_{i-1}\le2v_i+2d<4v_i$ by \eqref{eq:app-h-star}, so $q_i\le3$.  The final
remainder $X'<v_{i^*}$ is a multiple of $d$, since all values are, and is
appended as a block.  As $i^*-i_0\le k_V-3$, there are at most
$1+3(k_V-3)<3k_V$ codes, each of length $\le\theta+2d\le\ell_z$, and
$X'<v_{i^*}\le B_0+2d\le(6k_V^2+4k_V)\theta\le2\ell_z$.
\end{proof}

\subsection{Proof of Theorem~\ref{thm:two-letter-fanout-one}}

Let $A$ have a pure yield.  By (I1) the yields of $A$ of type $s(\mathtt a)$
are its pure yields, so $\tau(A,s(\mathtt a))\le n$, and it suffices to show
$\tau(A,s(\mathtt a))\le11\cdot3^{M^*}k_V^2\theta$, where
\[
 M^*:=\max\bigl(\{5\}\cup\{M\ge6:(27/4)^{\lfloor M/3\rfloor}
 \le(M+1)^4k_V(3k_V+M+4)^2\ell_z^{\,4}\}\bigr),
\]
a finite number since the left side grows exponentially in $M$.  If the root
subtree has at most $k_V$ seeds, or $n_{X_{i^*}}=n$, then $n\le2k_V(k_V+2)\theta$.
Otherwise take the indices of Lemma~\ref{lem:app-h-select}.

\emph{Counting.}  Suppose $M\ge6$.  Let $X:=V^*-\sum_tw_t$; then $0<X<n$, since
$\sum_tw_t<\Lambda\le V^*\le n$ and $w_1>0$, and $X$ is a multiple of $d$.  Let
$Q$ be given by Lemma~\ref{lem:app-h-filler}.  The codes $\hat z_{j_1},\ldots,\hat
z_{j_M}$ have length $\le\ell_z$ and superincreasing values, so their subset
sums are distinct, and $\sum_tw_t<\Lambda\le D$.  With
$|l_\circ|+|r_\circ|\le\theta+2d\le\ell_z$ we get
$m\le(3k_V+M+3)\ell_z$, and Lemma~\ref{lem:app-h-count}, with
$(m+1)^2(M\ell_z+1)^2\le(3k_V+M+4)^2(M+1)^2\ell_z^{\,4}$, gives $M\le M^*$.  So in
every case $M\le M^*$, and by Lemma~\ref{lem:app-h-select}
\[
 \Lambda\le3^{M^*}(v_{i^*}+d)\le3^{M^*}(B_0+3d)\le10\cdot3^{M^*}k_V^2\theta .
\]

\emph{Conclusion.}  If $\Lambda=V^*$, then $n=V^*+\rho+\rho'<\Lambda+2d$ and
$\tau(A,s(\mathtt a))\le n$.  If $\Lambda=D<V^*$, then $P_0\ne\{0\}$, and by
Lemma~\ref{lem:app-h-period}(iii) $\mathtt a^n\equiv_{L,s}\mathtt a^{n'}$ with
$n'\le p_0\le D$, so $\tau(A,s(\mathtt a))\le n'$.  Together with the small cases,
$\tau(A,s(\mathtt a))\le\max\{2k_V(k_V+2)\theta,\,10\cdot3^{M^*}k_V^2\theta+4k_V^2\theta\}
\le11\cdot3^{M^*}k_V^2\theta$.

\emph{The exponents.}  Put $\mu:=(27/4)^{1/3}$, $\gamma:=\ln3/\ln\mu<1.7260$,
$\eta:=1/80$ and $C_\eta:=\sup_{M\ge0}(M+1)^6\mu^{-\eta M}<e^{34}$ (the maximum of
$6\ln x-(x-1)\eta\ln\mu$ over $x>0$ is attained at $x=6/(\eta\ln\mu)$).  For $M\ge6$
in the set defining $M^*$, $\mu^{M-2}\le(27/4)^{\lfloor M/3\rfloor}$ and
$3k_V+M+4\le5k_V(M+1)$ give
\[
 \mu^M\le15625\,\mu^2(M+1)^6k_V^{11}\theta^4,\qquad\text{hence}\qquad
 \mu^{(1-\eta)M}\le15625\,\mu^2C_\eta k_V^{11}\theta^4,
\]
and $3^M=\mu^{\gamma M}\le(15625\,\mu^2C_\eta k_V^{11}\theta^4)^{\gamma/(1-\eta)}$;
this also holds for $M=5$.  Since $\gamma/(1-\eta)=80\gamma/79<1.7479$, we obtain
\[
 \tau(A,s(\mathtt a))\le11\cdot3^{M^*}k_V^2\theta\le C_0k_V^{2+19.23}\theta^{1+6.992}
 \le C_0k_V^{22}\theta^8,\qquad C_0:=11(15625\,\mu^2C_\eta)^{80\gamma/79}<10^{36}.
\]
As $k_V\le2\|G\|$, $\theta\le\bar\theta_G$, and the nonterminals of $G$ are
nonterminals of $G'$ with the same yields, the theorem follows with
$C:=2^{22}C_0<10^{43}$.  Letting $\eta\to0$ shows that the exponents can be taken
arbitrarily close to $2+11\gamma<20.99$ and $1+4\gamma<7.91$.\qed

\subsection{Other Types and Proof of Corollary~\ref{cor:two-letter-fanout-one}}

\begin{lemma}
\label{lem:app-h-other}
Let $G'$ be as above and let $A$ be a nonterminal of $G'$.
\begin{enumerate}[label=\textup{(\roman*)},leftmargin=*]
\item If $s:\Sigma^+\twoheadrightarrow U_1$ and $\mathtt a$ is the only letter of
type $1$, and $A$ has a yield of type $\zeta$, then
$\tau(A,\zeta)\le(k_V-1)\theta+1$.
\item If $\Sigma=\{\mathtt a,\mathtt b\}$, $s=c$, and $A$ has a yield of content
$\Sigma$, then $\tau(A,\Sigma)\le3k_V\theta+2$.
\end{enumerate}
\end{lemma}

\begin{proof}
(i) Take a derivation of a yield of $A$ of type $\zeta$ with the least number of
nodes and mark a leaf whose letter is not $\mathtt a$.  Replacing the subtree of a
node on the path from the root to the marked leaf by the subtree of a lower node
on the same path with the same label keeps the marked leaf, hence the type; so no label repeats on
this path, which has at most $k_V$ nonterminal nodes.  Replacing the at most
$k_V-1$ subtrees hanging off it by yields of length $\le\theta$ gives a yield of
$A$ of type $\zeta$ and length $\le(k_V-1)\theta+1$.  (ii) Take a derivation of a
yield of content $\Sigma$ with the least number of nodes and mark an $\mathtt
a$-leaf and a $\mathtt b$-leaf.  By the same contraction within each of the three
path segments between the root, the branching node and the two leaves (never
between different segments), the skeleton has at most
$3k_V$ nodes; replacing the at most $3k_V$ subtrees hanging off it by yields of
length $\le\theta$ gives a yield of content $\Sigma$ of length $\le3k_V\theta+2$.
In both cases the new word is a yield of $A$ of the same type, hence in the same
class (Lemma~\ref{lem:typed-homogeneity}).
\end{proof}

\begin{proof}[Proof of Corollary~\ref{cor:two-letter-fanout-one}]
In both cases $G'$ is a reduced good presentation of $L\cap\Sigma^+$ of fan-out
one and rank $2\le r_{\max}$, with $\bar\theta_{G'}\le\bar\theta_G$, and
Corollary~\ref{cor:distribution-representatives} applied to $G'$ gives
$\|K\|_+\le N_{\mathrm{typ}}(G',s)(1+2\bar\theta_G+r_{\max}\tau(G',s))
+\mathbf 1_{\{\eps\in L\}}$.  It remains to bound $\tau(G',s)$ and
$N_{\mathrm{typ}}(G',s)$.

(i) The positive image of $s$ is generated by $s(\mathtt a)$ and $s(\mathtt b)$ in a
semilattice, so $s(w)$ depends only on $c(w)$; thus $s\preceq c$,
$\Cmcf{1}{s}\subseteq\Cmcf{1}{c}$, and $\equiv_{L,c}$ refines $\equiv_{L,s}$.
For a typed nonterminal $A_t$ of $G'$ and the $c$-type $t'$ of one of its
yields, $\tau_{A_t}(G',s)\le\tau(A,t')$.  Both letters are isolated for $c$, so
Theorem~\ref{thm:two-letter-fanout-one}, applied to $c$ with $\mathtt a$ and with
$\mathtt b$, and Lemma~\ref{lem:app-h-other}(ii) give
$\tau(G',s)\le C(\|G\|+2)^{22}\bar\theta_G^{\,8}$.  The typed refinement of $G'$
has at most $3k_V$ typed nonterminals and at most nine typed rules per rule, so
$N_{\mathrm{typ}}(G',s)\le3k_V+9|P'|\le12(\|G\|+2)$.

(ii) The typed nonterminals of $G'$ are of the form $A_1$, with $A$ having a
yield in $\mathtt a^+$, and $A_\zeta$.  By
Theorem~\ref{thm:two-letter-fanout-one} and Lemma~\ref{lem:app-h-other}(i),
$\tau(G',s)\le C(\|G\|+2)^{22}\bar\theta_G^{\,8}$.  There are at most $2k_V$ typed
nonterminals and at most four typed rules per rule, so
$N_{\mathrm{typ}}(G',s)\le2k_V+4|P'|\le6(\|G\|+\sigma_G)\le12\|G\|$.

In both cases $1+2\bar\theta_G\le3\bar\theta_G$ gives the bound with
$C':=12(C+2)<10^{45}$.
\end{proof}

\section{Typings onto \texorpdfstring{$U_1$}{U1} at Fan-Out One}
\label{app:u1-typings}

\subsection{Letter-Disjoint Images of Finite Substitutable Languages}

\begin{lemma}\label{lem:app-u1-insertion}
Let $L\subseteq\Delta^+$ be finite and Clark--Eyraud substitutable.  There are no
$P,Q\in\Delta^*$ and $T\in\Delta^+$ with $PQ\in L$ and $PTQ\in L$.
\end{lemma}

\begin{proof}
Suppose otherwise.  If $P\ne\eps$, the words $P$ and $PT$ share the context
$(\eps,Q)$, so $\D_L(P)=\D_L(PT)$; then $PT^jQ\in L$ implies $PT^{j+1}Q\in L$, and
$L$ contains the words $PT^jQ$ of every length, contrary to finiteness.  If
$P=\eps$, argue in the same way with $Q$ and $TQ$, which share the context
$(\eps,\eps)$.
\end{proof}

\begin{lemma}\label{lem:app-u1-code}
Let $h:\Delta^*\to\Sigma^*$ be a morphism such that every $h(\delta)$,
$\delta\in\Delta$, is nonempty and no letter occurs twice in the concatenation of
all $h(\delta)$.  If $L\subseteq\Delta^+$ is finite and Clark--Eyraud
substitutable, so is $h(L)$.
\end{lemma}

\begin{proof}
Each letter occurring in some $h(\delta)$ determines $\delta$ and its position in
$h(\delta)$.  Hence $h$ is injective, and in a word $h(D)$ every letter occurrence
belongs to a determined occurrence of a codeword $h(\delta)$; whether two adjacent
letters lie in one codeword occurrence, and whether a letter begins or ends one,
is read off the letters themselves.

Let $x,y\in\Sigma^+$ with $pxq=h(D_1)$ and $pyq=h(D_2)$ in $h(L)$, $x\ne y$.  The
codeword occurrences lying inside $p$, and those inside $q$, are the same in both
words; so $D_i=PM_iQ$, where $M_i$ consists of the codeword occurrences meeting
$x$, respectively $y$.  If $p$ ends inside a codeword occurrence $h(\delta)$, then
$M_1$ and $M_2$ both begin with $\delta$, and if $q$ begins inside one,
$h(\delta')$, both end with $\delta'$.

Suppose that $x$ lies inside one codeword occurrence $h(\delta)$ and is not all
of it; say $p$ ends inside it (the other case is symmetric).  Then $M_1=\delta$
and $M_2$ begins with $\delta$; $M_2\ne M_1$, since otherwise $y=x$, so
$M_2=\delta T$ with $T\in\Delta^+$, and $D_1=P\delta Q$, $D_2=P\delta TQ$ contradict
Lemma~\ref{lem:app-u1-insertion}.  The
same holds with $x$ and $y$ exchanged.  So neither lies inside one codeword
occurrence unless it is a whole codeword.  Then $M_1=[\delta]X[\delta']$ and
$M_2=[\delta]Y[\delta']$, where $X,Y\in\Delta^*$ are the codewords lying inside $x$,
respectively $y$, and the bracketed letters are present iff the corresponding
cut lies inside a codeword.  If $X=\eps\ne Y$ or $Y=\eps\ne X$,
Lemma~\ref{lem:app-u1-insertion} is contradicted again, and $X=Y=\eps$ gives
$x=y$.  Hence $X,Y\in\Delta^+$, and
\[
 x=\alpha\,h(X)\,\beta,\qquad y=\alpha\,h(Y)\,\beta,
\]
where $\alpha$ is the part of $h(\delta)$ after $p$ (or $\eps$) and $\beta$ the
part of $h(\delta')$ before $q$ (or $\eps$).  The words $X$ and $Y$ share a context
in $L$, so $\D_L(X)=\D_L(Y)$.

Now let $p'xq'\in h(L)$.  The cut between $\alpha$ and $h(X)$, and the one between
$h(X)$ and $\beta$, are boundaries of codeword occurrences, which is read off the
letters of $x$; so $p'\alpha=h(P')$ and $\beta q'=h(Q')$ for some $P',Q'$, and
$P'XQ'\in L$.  Then $P'YQ'\in L$, and $h(P'YQ')=p'yq'$.  By symmetry
$\D_{h(L)}(x)=\D_{h(L)}(y)$.
\end{proof}


\subsection{Quadratic Data for the Binary Typing}

We prove Proposition~\ref{prop:binary-u1-quadratic}.  Put
$k:=2^m$, $M:=2^n$ and $\psi(w):=|w|_{\mathtt b}+|w|_{\mathtt a}/k$, which is
additive.  For a word $w$ let $\lambda(w)$ and $\varrho(w)$ be the lengths of its
longest suffix and longest prefix in $\mathtt a^*$.

\begin{proof}[Proof of Proposition~\ref{prop:binary-u1-quadratic}]
(i) A word of $\Sigma^+$ lies in $L_{m,n}$ iff all its maximal runs of $\mathtt a$
have length divisible by $k$ and $\psi(w)=M$: such a word factors uniquely into
blocks $\mathtt a^k$ and $\mathtt b$, and $\psi$ counts the blocks.  Let
$x,y\in\Sigma^+$ have the same type and let $uxv,uyv\in L_{m,n}$.  By additivity
$\psi(x)=\psi(y)$.  If $x,y$ are pure, this gives $x=y$.  Otherwise write
$x=\mathtt a^{\alpha}x'\mathtt a^{\beta}$ with $x'$ beginning and ending with
$\mathtt b$, and similarly $y=\mathtt a^{\alpha'}y'\mathtt a^{\beta'}$.  The maximal
$\mathtt a$-runs inside $x'$ and $y'$ are maximal runs of $uxv$ and $uyv$, so their
lengths are divisible by $k$; the runs through the two ends of $x$ have lengths
$\lambda(u)+\alpha$ and $\beta+\varrho(v)$, so $\alpha\equiv\alpha'$ and
$\beta\equiv\beta'\pmod k$.  For any context $(p,q)$ the maximal $\mathtt a$-runs of
$pxq$ are those of $p$ not at its end, a run of length $\lambda(p)+\alpha$, the
runs inside $x'$, a run of length $\beta+\varrho(q)$, and those of $q$ not at its
start; the same holds for $pyq$ with $\alpha',\beta',y'$.  So the divisibility
conditions agree, $\psi(pxq)=\psi(pyq)$, and $pxq\in L_{m,n}\iff pyq\in L_{m,n}$.
Finite languages are context-free.

(ii) Take nonterminals $C_0,\ldots,C_m$, $A_0,\ldots,A_n$, start symbol $A_n$, and the
rules $C_0\to\mathtt a$, $C_j\to C_{j-1}C_{j-1}$ ($1\le j\le m$),
$A_0\to C_m\mid\mathtt b$ and $A_i\to A_{i-1}A_{i-1}$ ($1\le i\le n$).  Then
$L(C_j)=\{\mathtt a^{2^j}\}$, $L(A_0)=\{\mathtt a^k,\mathtt b\}$ and $L(A_i)$ is the
set of concatenations of $2^i$ words of $L(A_0)$, so $L(A_n)=L_{m,n}$; the grammar
is good and reduced.  It has $m+n+2$ nonterminals and rules of sizes $2$, $5$
($m$ rules), $3$, $2$ and $5$ ($n$ rules) (Definition~\ref{def:standard-mcfg}),
which gives $6m+6n+9$.  At fan-out one, the shortest sentence using a rule
$R:B\to\rho(B_1,\ldots,B_q)$ has length $\theta_G(R)=o_B+c_R+\sum_jd_{B_j}$, where
$d_C$ is the length of a shortest yield of $C$, $o_C$ the least number of letters
around an occurrence of $C$ in a sentence, and $c_R$ the number of terminals
written by $R$; here
$d_{A_0}=1$, $d_{A_i}=2^i$, $d_{C_j}=2^j$, $o_{A_i}=M-2^i$, $o_{A_0}=o_{C_m}=M-1$ and
$o_{C_j}=M-1+k-2^j$.  This gives $\theta_G(R)=M$ for the rules of $A_i$ and for
$A_0\to\mathtt b$, and $\theta_G(R)=M+k-1$ for $A_0\to C_m$ and the rules of the
$C_j$; so $\bar\theta_{G_{m,n}}=2^m+2^n-1$.

(iii) The only pure word of $L_{m,n}$ is $\mathtt a^{kM}$, so by
Proposition~\ref{prop:rare-type} every characteristic sample contains it and
$\|F\|_+\ge kM+1>2^{m+n}$.  For $m=n$, $\bar\theta_{G_{n,n}}<2^{n+1}$ gives
$\bar\theta_{G_{n,n}}^{\,2}/4<4^n<\|F\|_+$.  A bound $p(\|G\|)\bar\theta_G$ would give
$4^n\le p(12n+9)\,2^{n+1}$ for all $n$, which fails for every polynomial $p$.
\end{proof}


\section{Content Typings on Two Letters at Fan-Out at Least Two}
\label{app:two-letter-content}

This appendix proves Theorem~\ref{thm:two-letter-content}.  Throughout,
$\Sigma=\{\mathtt a,\mathtt b\}$, $c:=c_\Sigma$, $f\ge2$ is fixed, $L\in\Cmcf{f}{c}$, and
$G=(V,\Sigma,P,S,\fan)$ is a reduced good grammar of fan-out $\le f$ and rank
$\le r:=r_{\max}$ with $L(G)=L\cap\Sigma^+$; put $v:=|V|$ and
$\theta:=\bar\theta_G$; we also put $D_f:=(6f+1)3^f$.  Apart from Lemma~\ref{lem:app-two-basic}(i), which uses the
condition of the arity of the typed nonterminal at hand, the arguments use only
$(1,c)$- and $(2,c)$-tuple-substitutability, which hold since $f\ge2$.  For each $A\in V$ the context $E_A$ is that of
Lemma~\ref{lem:thickness-equivalence}(iv), so $|E_A|_{\mathrm{fix}}\le\theta$
and $E_A[L_A(G)]\subseteq L(G)$; for each $B\in V$ fix $\omega_B\in L_B(G)$ with
$\|\omega_B\|_1\le\theta$ (item~(ii)), and recall that every rule writes at most
$\theta$ terminals (item~(i)).  In a derivation tree, the \emph{yield} of a node
$u$ is the tuple it derives, $\ell_u$ is its total length, $B_u$ its
nonterminal, and $u$ is \emph{heavy} if $\ell_u>\theta$; a node \emph{writes}
the terminals of its template.  The letters $x,y$ always denote the two
letters of $\Sigma$ in some order.  A \emph{run} of a word is a maximal factor
$x^k$ ($k\ge1$); a run is \emph{internal} if it is neither the first nor the
last run, so that an internal $y$-run is preceded and followed by $x$.  The
\emph{$x$-skeleton} $\sigma_x(w)$ replaces every $x$-run of $w$ by one $x$, and
$w$ has a \emph{spare} $x$ if $|w|_x$ exceeds the number of $x$-runs of $w$.
We say that \emph{$x^k$ occurs} if some word of $L$ has the factor $x^k$, and
that $L$ is \emph{Parikh-saturated} if for $u,w\in\Sigma^+$ with the same
Parikh vector $\Psi(u)=\Psi(w)$ we have $u\in L$ iff $w\in L$.

\paragraph{Map of the proof.}
The table lists, for each step, what it uses and where it is used.
\begin{center}
\footnotesize
\setlength{\tabcolsep}{3pt}
\renewcommand{\arraystretch}{1.15}
\begin{tabularx}{\linewidth}{@{}>{\raggedright\arraybackslash}p{0.24\linewidth}>{\raggedright\arraybackslash}X>{\raggedright\arraybackslash}p{0.26\linewidth}@{}}
\hline
Step & Uses & Used in\\
\hline
Lemma~\ref{lem:app-two-basic}(i) typed closure
 & Lemma~\ref{lem:typed-homogeneity}, arity-$d$ substitutability
 & Lemma~\ref{lem:app-two-saturated}\\
Lemma~\ref{lem:app-two-basic}(ii) cutting identity paths
 & goodness of $G$
 & (P2), the mixed case of Lemma~\ref{lem:app-two-saturated}\\
Lemma~\ref{lem:app-two-runs} run exchanges
 & $(1,c)$, $(2,c)$-substitutability, $x^4$ occurs
 & Lemma~\ref{lem:app-two-alternate}, Proposition~\ref{prop:app-two-nonsaturated}\\
Lemma~\ref{lem:app-two-alternate}
 & Lemma~\ref{lem:app-two-runs}(i),(iii)
 & Proposition~\ref{prop:app-two-nonsaturated}\\
Lemma~\ref{lem:app-two-insertion} runs created by replacement
 & Lemma~\ref{lem:witness-length-template}(i), node-minimal shortest $T$
 & Proposition~\ref{prop:app-two-nonsaturated}\\
Proposition~\ref{prop:app-two-nonsaturated} (not saturated)
 & the three lemmas above; typed labels on paths
 & Theorem~\ref{thm:two-letter-content}\\
Lemma~\ref{lem:app-two-saturated}, union $\Sigma$
 & Lemma~\ref{lem:app-two-basic}; floors along paths
 & Theorem~\ref{thm:two-letter-content}\\
(P1) row rigidity
 & $(1,c)$-substitutability, minimality of $K$
 & (P3), (P4)\\
(P2) height, mixed nodes
 & Lemma~\ref{lem:app-two-basic}(ii), minimality of $K$
 & (P3), (P6)\\
(P3) heavy pure nodes
 & (P1), (P2)
 & (P4), (P5)\\
(P4) common ratio $\alpha$
 & (P3)
 & (P5), (P6)\\
(P5) bound on $\alpha$
 & (P3), (P4)
 & (P6)\\
(P6) counting
 & (P2), (P4), (P5)
 & Theorem~\ref{thm:two-letter-content}\\
\hline
\end{tabularx}
\end{center}

\begin{lemma}\label{lem:app-two-basic}
\textup{(i)} For every typed nonterminal $X=A_{\mathbf t}$ of arity $d$,
$\operatorname{Cl}_{L,c}(X)=\{\tuple z\in(\Sigma^+)^d:c^{(d)}(\tuple z)=\mathbf t,\
E_A[\tuple z]\in L\}$; hence $\tau_X$ is the least total length of such a tuple.

\textup{(ii)} In a derivation tree of $G$, let $u$ be a proper ancestor of $u'$
with the same nonterminal, and suppose their yields have the same total Parikh
vector.  Then every rule on the path from $u$ to $u'$ has rank one, writes no
terminal and is a componentwise identity, so replacing the subtree at $u$ by the
subtree at $u'$ changes no derived tuple and removes at least one node.
\end{lemma}

\begin{proof}
(i) Fix $\tuple u\in L_X$.  By Lemma~\ref{lem:typed-homogeneity},
$\operatorname{Cl}_{L,c}(X)$ is the set of tuples with the type and the
distribution of $\tuple u$; they are accepted by $E_A\in\D_L(\tuple u)$.  A tuple
$\tuple z$ of type $\mathbf t$ accepted by $E_A$ shares this context with $\tuple u$,
so $\D_L(\tuple z)=\D_L(\tuple u)$ by substitutability.

(ii) By linearity and nondeletion every letter of the yield of $u'$ occurs in the
yield of $u$, so equal Parikh vectors exclude terminals written on the path
and siblings off the path, since siblings have nonempty yields
($\eps$-freeness).  So every rule on the path has one child of some fan-out
$d'$, each of its $d'$ variables occurs once (nondeletion) and in order
(nonpermutation), and each parent component is nonempty ($\eps$-freeness).
Two variables in one component would be consecutive, which nonmerging forbids;
hence every component is a single variable, the fan-outs agree, and the rule is
the identity.
\end{proof}

\begin{lemma}[Run exchanges]\label{lem:app-two-runs}
Suppose that $x^4$ occurs.
\begin{enumerate}[label=\textup{(\roman*)},leftmargin=*]
\item If $w,w'\in\Sigma^+$ satisfy $\sigma_x(w)=\sigma_x(w')$ and
$|w|_x=|w'|_x$, then $w\in L$ iff $w'\in L$.
\item If some word of $L$ has a factor $xxyx$ or $xyxx$, then $L$ is
Parikh-saturated.
\item If $L$ is not Parikh-saturated, then in every word of $L$ with a factor
$xx$ every internal $y$-run has length at least $2$.
\item If $L$ is not Parikh-saturated and $y^4$ does not occur, then every word
of $L$ with a factor $xx$ has at most one internal $y$-run.
\end{enumerate}
\end{lemma}

\begin{proof}
(i) Since $c(x)$ is idempotent, Lemma~\ref{lem:power-transfer} applies with
$u=x$, $i_0=p=1$ and $i=\delta=j=1$: as $x^4$ occurs, $w_{\mathrm L}xw_{\mathrm M}xxw_{\mathrm R}\in L$ iff
$w_{\mathrm L}xxw_{\mathrm M}xw_{\mathrm R}\in L$ whenever $w_{\mathrm M}\ne\eps$.  Let $R$ and $R'$ be $x$-runs of a word $w$ with
$|R|\ge2$, say $R'$ to the left of $R$ (the other case is symmetric).  Write
$w=U\,x\,Z\,xx\,V$, where the first displayed $x$ is the last letter of $R'$ and
$xx$ are the first two letters of $R$; the factor $Z$ separates two runs, so
$Z\ne\eps$.  Hence $U\,xx\,Z\,x\,V\in L$ iff $w\in L$.  This moves one $x$ from $R$
to $R'$ and keeps the $x$-skeleton and $|w|_x$.  Words with the same
$x$-skeleton have the same number $\rho$ of $x$-runs, and their run lengths are
compositions of $|w|_x$ into $\rho$ positive parts; any two such compositions
are connected by moves of one unit from a part of size at least $2$ to another
part.

(ii) Say $UxxyxV\in L$ (the other case is symmetric).  By (i),
$UxyxxV\in L$.  The context $(Ux,xV)$ accepts $xy$ and $yx$, both of content
$\Sigma$, so $\D_L(xy)=\D_L(yx)$: in every context $xy$ may be replaced by $yx$
and conversely.  Adjacent transpositions connect all words with one Parikh
vector.

(iii) Let $w\in L$ have a factor $xx$ and an internal $y$-run of length $1$.
Then $w$ has a spare $x$, and by (i) the spare letters can be moved into the
$x$-run preceding that $y$-run; the resulting word of $L$ has a factor $xxyx$,
so $L$ is Parikh-saturated by (ii).

(iv) Suppose that $w\in L$ has a factor $xx$ and two internal $y$-runs.  Then it
has two consecutive internal $y$-runs $y^e$ and $y^{e'}$, separated by an
$x$-run $R$; let $A_0$ be the $x$-run preceding $y^e$.  By (iii), $e,e'\ge2$.
Since $w$ has a spare $x$, (i) gives two words of $L$ with the $x$-skeleton and
the number of letters $x$ of $w$: a word $W_2=U\,x\,y^e\,xx\,V'$ in which $R$ has
length $2$ (the displayed $x$ before $y^e$ is the last letter of $A_0$, and $V'$
begins with $y^{e'}$), and a word $W_1=U_1\,x\,y^e\,x\,y^{e'}V_1$ in which $R$ has
length $1$.  The context $(U,xV')$ accepts $xy^ex$, giving $W_2$, and $xxy^e$,
giving $U\,xx\,y^e\,x\,V'$, which has the $x$-skeleton and the $x$-count of $W_2$
and lies in $L$ by (i).  Both tuples have content $\Sigma$, so
$\D_L(xy^ex)=\D_L(xxy^e)$.  The context $(U_1,y^{e'}V_1)$ accepts $xy^ex$, giving
$W_1$, hence also $xxy^e$, so $U_1\,xx\,y^{e+e'}V_1\in L$.  As $e+e'\ge4$, this
word has the factor $y^4$, a contradiction.
\end{proof}

\begin{lemma}\label{lem:app-two-alternate}
If $L$ is not Parikh-saturated and both $\mathtt a^4$ and $\mathtt b^4$ occur,
then every run of every word of $L$ with at least four runs has length $1$.
\end{lemma}

\begin{proof}
Let $w\in L$ have at least four runs.  Its second and third runs are internal
and consist of different letters, so for each letter $y$ the word $w$ has an
internal $y$-run $\beta$ and at least two $y$-runs.  By
Lemma~\ref{lem:app-two-runs}(i) for $y$, moving letters $y$ from $\beta$ to another
$y$-run gives a word of $L$ in which $\beta$ has length $1$ and the $x$-runs are
those of $w$.  If an $x$-run of $w$ had length at least $2$, this word would
have a factor $xx$ and an internal $y$-run of length $1$, contrary to
Lemma~\ref{lem:app-two-runs}(iii).  So all $x$-runs of $w$ have length $1$, for
both choices of $x$.
\end{proof}

\begin{lemma}[Runs created by short replacements]\label{lem:app-two-insertion}
Let $X=A_{\mathbf t}$ be a typed nonterminal of arity $d\le f$, and let $T$ be a
derivation of $\widetilde G_0$ from $X$ whose yield $\tuple u$ has the least total
length and, among those, the fewest nodes.  If $T$ has an antichain $Q$ of
heavy nodes with $|Q|\ge6f+2$, then for some letter $x$:
\textup{(a)} some component of $\tuple u$ has content $\{x\}$ and length at least
$5$, so $x^5$ occurs; and
\textup{(b)} some word $W\in L$ has a factor $xx$ and at least four runs, and,
unless both $\mathtt a^5$ and $\mathtt b^5$ occur, $W$ has two internal $y$-runs.
\end{lemma}

\begin{proof}
In each component $u_i$ of content $\{x\}$ (a \emph{pure} component) mark the
first and the last $\min\{2,|u_i|\}$ letters, and in each component of content
$\Sigma$ mark one $\mathtt a$ and one $\mathtt b$; there are at most $4f$ marks.
Each mark is written by one node, and nodes of $Q$ have disjoint subtrees, so
the set $Q_0$ of nodes of $Q$ whose subtrees write no mark has
$|Q_0|\ge|Q|-4f\ge2f+2$.

Let $h\in Q_0$ with nonterminal $C_h$.  Erase the type indices and replace the
subtree at $h$ by a derivation of $\omega_{C_h}$.  The new root tuple lies in
$L_A(G)$ and is shorter, since by nondeletion every letter of the yield of $h$
occurs in $\tuple u$ and $\|\omega_{C_h}\|_1\le\theta<\ell_h$.  The marks keep the
content of every component of content $\Sigma$.  If every pure component stayed
pure, the new tuple would have type $\mathbf t$ and would lie in $L_X$ by
Lemma~\ref{lem:witness-length-template}(i) with $\lambda:=\widehat c$, contrary to
the choice of $T$.  Hence $\omega_{C_h}$ puts the other letter into a block of $h$
(an interval of a root component occupied by a component of the yield of $h$)
that lies in a pure component.  Blocks contain no mark, so that pure component
has length at least $5$, and the block lies after its first two and before its
last two letters.  This proves (a), since $E_A[\tuple u]\in L$.

Each $h\in Q_0$ owns at most $f$ blocks.  Join two owners when blocks of theirs in
one root component are adjacent with no letter between them.  Each block has at
most two adjacent blocks, so this graph has maximum degree at most $2f$ and an
independent set $I$ with $|I|\ge|Q_0|/(2f+1)>1$, hence $|I|\ge2$.  Replace the subtrees at
all $h\in I$ at once, giving $\tuple u'\in L_A(G)$ and $W:=E_A[\tuple u']\in L$.
Fix a pure component of letter $x$ and write $y$ for the other letter.  Its
letters outside the blocks of $I$ are unchanged, hence equal to $x$; this
covers the blocks of owners outside $I$, which are kept, and the first two and
the last two letters, which are marked.  So every maximal $y$-run of the new
component is preceded and followed by $x$, and it lies in a union of blocks of
$I$ that are consecutive with no letter between them.  Such blocks belong to one
owner: two blocks of distinct owners $h\ne h'$ of $I$ with no block of $I$
between them are separated by an unchanged letter, since their gap is nonempty
(otherwise $h$ and $h'$ would be joined) or contains a block of an owner outside
$I$.  Hence every $y$-run of the new component has a unique \emph{creating}
owner in $I$, and by the previous paragraph every $h\in I$ creates at least one
run, in a pure component whose letter we call $x_h$.  For $h\ne h'$ in $I$ the
runs they create are distinct: if they lie in the same pure component, they have
different creating owners; if they lie in different pure components, they are
in different components of $\tuple u'$.

Take $h_1\ne h_2$ in $I$.  If $x_{h_1}=x_{h_2}=:x$, then $W$ has the factor $xx$
(the first two letters of a pure component of length at least $5$) and two
internal $y$-runs, hence at least five runs.  Otherwise both $\mathtt a^5$ and
$\mathtt b^5$ occur, by the argument for (a) applied to $h_1$ and to $h_2$, and with $x:=x_{h_1}$ and $y:=x_{h_2}$ the two pure
components contain the runs $x,y,x$ and $y,x,y$ respectively; these are six
distinct runs of $W$, and $W$ has the factor $xx$.
\end{proof}

\begin{proposition}\label{prop:app-two-nonsaturated}
If $L$ is not Parikh-saturated, then $t_{\mathrm{typ}}(G,c)\le D_f(r+1)v\theta$.
\end{proposition}

\begin{proof}
Let $X$ and $T$ be as in Lemma~\ref{lem:app-two-insertion}.  We first show that
every antichain of heavy nodes of $T$ has at most $6f+1$ nodes.  Otherwise
Lemma~\ref{lem:app-two-insertion} applies.  If neither $\mathtt a^4$ nor
$\mathtt b^4$ occurs, (a) is impossible.  If both occur, the word $W$ of (b)
has at least four runs and a run of length at least $2$, contrary to
Lemma~\ref{lem:app-two-alternate}.  If exactly one occurs, say $x_0^4$, then by
(a) $x=x_0$ and $y^5$ does not occur, so by (b) $W$ has a factor $xx$ and two
internal $y$-runs, contrary to Lemma~\ref{lem:app-two-runs}(iv).

A typed label is a nonterminal together with the contents of its components,
so there are at most $3^fv$ of them.  No typed label repeats on a path of $T$,
since cutting between two equal labels gives a derivation from $X$ whose yield
is not longer and which has fewer nodes; so paths have at most $3^fv$ nodes.
Heavy nodes are closed under ancestors, and those at one depth form an
antichain, so there are at most $(6f+1)3^fv=D_fv$ of them.  If the root is not
heavy, $\|\tuple u\|_1\le\theta$.  Otherwise every letter of $\tuple u$ is written by
a heavy node or lies in the yield of a non-heavy child of one, and each heavy
node writes at most $\theta$ letters and has at most $r$ children.  Hence
$\|\tuple u\|_1\le D_f(r+1)v\theta$.
\end{proof}

\begin{lemma}\label{lem:app-two-saturated}
If $L$ is Parikh-saturated, then
$\tau(G,c)\le C_{f,r}v^2\theta^4H_\theta$, where $C_{f,r}:=(r+1)(4r(3f+1)+D_f)+3$.
\end{lemma}

\begin{proof}
Fix $X=A_{\mathbf t}$ of arity $d$, put $(p,q):=\Psi(E_A)$, so $p+q\le\theta$, and
$\mathcal P:=\Psi(L\cap\Sigma^+)$.  By Lemma~\ref{lem:app-two-basic}(i) and Parikh
saturation, $\tau_X$ is the least $\|\tuple z\|_1$ over tuples $\tuple z$ of type
$\mathbf t$ with $(p,q)+\sum_i\Psi(z_i)\in\mathcal P$.  Write $(x_B,y_B):=\Psi(\omega_B)$,
and for a node $u$ of a derivation tree let $(a_u,b_u)$ be the Parikh vector of
its yield.  If $u$ lies in a successful derivation of a word $w$, replacing its
subtree by a derivation of $\omega_{B_u}$ gives a word of $L$ with Parikh vector
$\Psi(w)-(a_u,b_u)+(x_{B_u},y_{B_u})$; for a proper descendant $u'$ with
$B_{u'}=B_u$, replacing the subtree at $u$ by that at $u'$ gives one with
$\Psi(w)-(a_u-a_{u'},b_u-b_{u'})$.  Every comparison below is of this kind: it
concerns Parikh vectors of words of $L$ yielded by complete successful
derivations obtained from a fixed one by such replacements, and by Parikh
saturation membership in $L$ of a word depends only on its Parikh vector.  The union of the contents $t_1,\dots,t_d$ is
$\Sigma$, or all of them equal one singleton.

\emph{The union is $\Sigma$.}  Let $t_{\mathtt a}$ and $t_{\mathtt b}$ be the
numbers of components whose content contains $\mathtt a$, resp.\ $\mathtt b$; both
lie in $\{1,\dots,f\}$.  A tuple of type $\mathbf t$ with $n_{\mathtt a}$ letters $\mathtt a$
and $n_{\mathtt b}$ letters $\mathtt b$ exists iff $n_{\mathtt a}\ge t_{\mathtt a}$ and
$n_{\mathtt b}\ge t_{\mathtt b}$, so $\tau_X$ is the least
$n_{\mathtt a}+n_{\mathtt b}$ with these bounds and
$(p+n_{\mathtt a},q+n_{\mathtt b})\in\mathcal P$.  Fix a minimizing pair, put
$c_{\mathtt a}:=n_{\mathtt a}-t_{\mathtt a}+1\ge1$ and
$c_{\mathtt b}:=n_{\mathtt b}-t_{\mathtt b}+1\ge1$, and fix a successful derivation
$T$ of $W:=\mathtt a^{p+n_{\mathtt a}}\mathtt b^{q+n_{\mathtt b}}$ with the fewest
nodes.  If a heavy node $u$ had $a_u<c_{\mathtt a}$ and $b_u<c_{\mathtt b}$, its
replacement would give a Parikh vector $(p+n'_{\mathtt a},q+n'_{\mathtt b})\in\mathcal
P$ with $n'_{\mathtt a}\ge t_{\mathtt a}$, $n'_{\mathtt b}\ge t_{\mathtt b}$ and
$n'_{\mathtt a}+n'_{\mathtt b}<n_{\mathtt a}+n_{\mathtt b}$; so every heavy node has
$a_u\ge c_{\mathtt a}$ or $b_u\ge c_{\mathtt b}$.  If $u'$ is a proper descendant of
$u$ with the same nonterminal, then $a_u-a_{u'}\ge c_{\mathtt a}$ or
$b_u-b_{u'}\ge c_{\mathtt b}$: a zero difference contradicts node-minimality
(Lemma~\ref{lem:app-two-basic}(ii)), and another small difference contradicts
the choice of the pair.  Along a path, $(\lfloor a_u/c_{\mathtt a}\rfloor,
\lfloor b_u/c_{\mathtt b}\rfloor)$ is componentwise nonincreasing, and no
nonterminal repeats while it is constant.  Since
$\lfloor(p+n)/(n-t+1)\rfloor\le p+t$ for $n\ge t\ge1$ and $p\ge0$ (as $(p+t)(n-t+1)-(p+n)=(n-t)(p+t-1)\ge0$), it
takes at most $1+(p+t_{\mathtt a})+(q+t_{\mathtt b})\le\theta+2f+1$ values, so paths
have at most $v(\theta+2f+1)$ nodes.  Heavy nodes at one depth have disjoint yields,
so there are at most $(p+n_{\mathtt a})/c_{\mathtt a}+(q+n_{\mathtt b})/c_{\mathtt
b}\le\theta+2f$ of them.  Heavy nodes are closed under ancestors, and every letter
is written by a heavy node or lies in one of its at most $r$ non-heavy
children; hence
$\tau_X\le|W|\le\max\{\theta,(r+1)\theta(\theta+2f)(\theta+2f+1)v\}\le(2f+1)(2f+2)(r+1)v\theta^3$.

\emph{All contents equal $\{\mathtt a\}$} (the case $\{\mathtt b\}$ is symmetric),
and $h:=d\in\{1,\dots,f\}$.  Let $K:=\min\{k\ge h:(p+k,q)\in\mathcal P\}$; then
$\tau_X=K$, and we may assume $K>\theta+2(f-1)\ge p+2(h-1)$, since otherwise
$\tau_X\le(2f-1)\theta$.  Fix a successful derivation
$T$ of $W:=\mathtt b^q\mathtt a^{p+K}$ with the fewest nodes.

(P1) \emph{Row rigidity}: there are no $(c_0,y),(c_0+\delta,y)\in\mathcal P$ with
$c_0\ge1$ and $1\le\delta\le K-h$.  Otherwise the context
$(\mathtt b^y\mathtt a^{c_0-1},\eps)$ accepts $\mathtt a$ and $\mathtt a^{1+\delta}$, so
$\D_L(\mathtt a)=\D_L(\mathtt a^{1+\delta})$ by $(1,c)$-substitutability (both have
content $\{\mathtt a\}$); the context $(\mathtt b^q\mathtt a^{p+K-\delta-1},\eps)$ accepts
$\mathtt a^{1+\delta}$, hence $\mathtt a$, and $(p+K-\delta,q)\in\mathcal P$ with
$h\le K-\delta<K$, contradicting the choice of $K$.

(P2) \emph{Height.}  If $u'$ is a proper descendant of $u$ with the same
nonterminal and $b_u=b_{u'}$, then $a_u-a_{u'}\ge K-h+1$: the value $0$
contradicts node-minimality (Lemma~\ref{lem:app-two-basic}(ii)), and a value in
$[1,K-h]$ contradicts the choice of $K$.  Going down a path, $b_u$ is
nonincreasing (by nondeletion the yield of a child occurs in that of its parent),
so the nodes of the path with one value of $b_u$ form a segment; $b_u$ takes at most $q+1$
values, and three nodes with one nonterminal and one value of $b_u$ would give
$p+K\ge a_u\ge2(K-h+1)$, false as $K>p+2h-2$.  So paths have at most $2v(q+1)$
nodes.  A node is \emph{mixed} if its yield contains $\mathtt b$.  For
$1\le k\le q$, the mixed nodes with $b_u=k$ that have no such proper ancestor have
pairwise disjoint yields, so there are at most $\lfloor q/k\rfloor$ of them; below
one of them, two such nodes are comparable (otherwise its yield would contain $2k$
letters $\mathtt b$), so they form a chain, which has at most $2v$ nodes by the
argument above.  Hence there are
$M\le2v\sum_{k=1}^q\lfloor q/k\rfloor\le2vqH_q\le2v\theta H_\theta$ mixed nodes.

(P3) \emph{Heavy pure nodes.}  Call a node $u$ that is not mixed a
\emph{heavy pure} node if $\ell_u=a_u>\theta+f-1$, and put
$g_u:=\ell_u-x_{B_u}\ge f$.  Two comparable heavy pure nodes have different
nonterminals, since by (P2) the upper one would have
$\ell_u\ge K-h+1+\theta+f>p+K=|W|_{\mathtt a}$.  In particular heavy pure nodes with
one nonterminal have disjoint yields.  For an antichain $J$ of heavy pure nodes,
replacing all of them gives a word with Parikh vector $(p+K-\sigma(J),q+Y(J))$,
where $\sigma(J):=\sum_{u\in J}g_u$ and $Y(J):=\sum_{u\in J}y_{B_u}$; this word is yielded by a complete successful derivation, so its Parikh vector
lies in $\mathcal P$.  Hence two antichains $J,J'$ with
$\sigma(J),\sigma(J')\le K-h$ and $Y(J)=Y(J')=:Y$ give
$(p+K-\sigma(J),q+Y),(p+K-\sigma(J'),q+Y)\in\mathcal P$, whose first coordinates
are at least $p+h\ge1$ and differ by $|\sigma(J)-\sigma(J')|\le K-h$; by (P1),
$\sigma(J)=\sigma(J')$.

(P4) \emph{Repeated nonterminals.}  If $B$ labels $m_B\ge2$ heavy pure nodes,
their yields are disjoint by (P3), so these nodes are pairwise incomparable and
every set of them is an antichain, so each has $\ell_u\le p+K-(\theta+f)\le K-f$ and
$g_u\le K-f\le K-h$; comparing two singletons by (P3) gives a common value $g_B$,
and then a common $\ell_B$.  If moreover $y_B=0$, one replacement gives
$(p+K-g_B,q)\in\mathcal P$ with $h\le K-g_B<K$, which is impossible; so $y_B\ge1$.
Call $B$ \emph{frequent} if $m_B>\theta$.  For frequent $B,C$, choose $y_C$ nodes
labelled $B$ and, separately, $y_B$ nodes labelled $C$, which is possible since
$m_B>\theta\ge y_C$ and $m_C>\theta\ge y_B$; both sets are antichains with
$Y=y_By_C$, and, since $y_C\le m_B-1$, $g_B\le\ell_B$ and the $m_B$ disjoint
yields give $m_B\ell_B\le p+K$, also $y_Cg_B\le(m_B-1)\ell_B\le p+K-\ell_B\le
K-f\le K-h$ (as $\ell_B\ge\theta+f\ge p+f$), and likewise for $C$.  By (P3), $y_Cg_B=y_Bg_C$, so all frequent $B$ share a ratio
$\alpha:=g_B/y_B$, and $\ell_B=x_B+\alpha y_B\le(1+\alpha)\theta$.

(P5) \emph{The ratio.}  Put $\alpha_*:=(r+1)(\theta+f-1)(2\theta+1)v\le3f(r+1)v\theta^2$;
we show $\alpha\le\alpha_*$ when frequent nonterminals exist.  This is clear if
$\alpha\le\theta$.  Otherwise fix a heavy pure node $u$ labelled by a frequent
nonterminal $B$, and let $H_u$ be the number of heavy pure nodes in its subtree.
These nodes are closed under ancestors within that subtree.  Every letter in the
yield of $u$ is written by a heavy pure node or belongs to the yield of a
non-heavy child of such a node.  Each heavy pure node writes at most $\theta$
letters and has at most $r$ such children, each of total yield length at most
$\theta+f-1$; hence $\ell_B\le(r+1)(\theta+f-1)H_u$.  Group these nodes by their
number of heavy pure proper ancestors in the subtree.  A chain of heavy pure
nodes has pairwise different nonterminals by (P3), hence at most $v$ nodes, so
this number is less than $v$ and there are at most $v$ groups; two comparable
nodes have different numbers, so each group is an antichain.  In one group the nodes with frequent labels have
length at least $\alpha$ and disjoint yields inside that of $u$, so there are at
most $(1+\alpha)\theta/\alpha=\theta+\theta/\alpha<\theta+1$, that is, at most
$\theta$, of them, so at most $v\theta$ over the at most $v$ groups; nodes with
other labels number at most $\theta$ per label, so at most $v\theta$ in all of
$T$.  So $H_u\le2v\theta\le v(2\theta+1)$, and
$\alpha\le g_B\le\ell_B\le\alpha_*$.

(P6) \emph{Counting.}  Call a node that is not mixed \emph{large} if
$\ell_u>T_*:=(\alpha_*+1)\theta\le(3f+1)(r+1)v\theta^3$.  By (P4)--(P5) no frequent
nonterminal labels a large node, so there are at most $v\theta$ large nodes, and
they are closed under ancestors among nodes that are not mixed.  Mixed nodes are closed under ancestors, since by
nondeletion the yield of a node occurs in that of its parent.  Let a letter
$\mathtt a$ of $W$ be written by a node $w$ that is neither mixed nor large.  By the
two closure properties, the ancestors of $w$ (including $w$) that are neither
mixed nor large form a final segment of the path from the root to $w$; let $z$ be
its top node.  Then $z$ is the root of $T$ or a child of a mixed or a large node,
and the subtree at $z$ contains no mixed or large node, so $\ell_z\le T_*$.  Thus
every letter $\mathtt a$ of $W$ is written by a mixed node, is written by a large
node, or lies in the yield of such a $z$; there are at most $1+rM+rv\theta$ nodes
$z$.  Hence
\[
\begin{aligned}
 p+K&\le\theta M+(rM+rv\theta+1)T_*+\theta\cdot v\theta\\
 &\le3v\theta^2H_\theta+4rv\theta H_\theta\cdot(3f+1)(r+1)v\theta^3
 \le\bigl(4r(3f+1)(r+1)+3\bigr)v^2\theta^4H_\theta,
\end{aligned}
\]
using $M\le2v\theta H_\theta$ and $rM+rv\theta+1\le4rv\theta H_\theta$.  As
$(2f+1)(2f+2)\le D_f$, in both cases, and when
$K\le\theta+2(f-1)$, $\tau_X\le C_{f,r}v^2\theta^4H_\theta$.
\end{proof}

\begin{proof}[Proof of Theorem~\ref{thm:two-letter-content}]
If $L$ is Parikh-saturated, Lemma~\ref{lem:app-two-saturated} applies.  Otherwise
Proposition~\ref{prop:app-two-nonsaturated} gives
$\tau(G,c)\le t_{\mathrm{typ}}(G,c)\le D_f(r+1)v\theta$.  In both cases
$\tau(G,c)\le C_{f,r}v^2\theta^4H_\theta$.  The image of $c$ has three elements, so
$|\mathsf S^\sharp|=4$ and $N_{\mathrm{typ}}\le C_c(G)\le4^{fr}\|G\|$.  By
Corollary~\ref{cor:distribution-representatives} there is a common characteristic
sample with
$\|K\|_+\le N_{\mathrm{typ}}(1+2\theta+r\tau(G,c))+\mathbf 1_{\{\eps\in L\}}$.  With
$v\le\|G\|$, this gives the first bound using $1+2\theta\le3\|G\|^2\theta^4H_\theta$ and
$\tau\le C_{f,r}\|G\|^2\theta^4H_\theta$, and the second using $1+2\theta\le3\|G\|\theta$ and
$\tau\le D_f(r+1)\|G\|\theta$.
\end{proof}
